% concatenated from main.tex
%%%%%%%%%%%%%%%%%%%%%%%%%%%%%%%%%%%%%%%%%%%%%%%%%%%%%%%%%%%%%%%%%%%%%%%%%%%%%%%%
%2345678901234567890123456789012345678901234567890123456789012345678901234567890
%        1         2         3         4         5         6         7         8
\documentclass[letterpaper, 10 pt, conference]{ieeeconf}  % Comment this line out if you need a4paper

%\documentclass[a4paper, 10pt, conference]{ieeeconf}      % Use this line for a4 paper

\IEEEoverridecommandlockouts                              % This command is only needed if 
                                                          % you want to use the \thanks command

% \overrideIEEEmargins                                      % Needed to meet printer requirements.

%In case you encounter the following error:
%Error 1010 The PDF file may be corrupt (unable to open PDF file) OR
%Error 1000 An error occurred while parsing a contents stream. Unable to analyze the PDF file.
%This is a known problem with pdfLaTeX conversion filter. The file cannot be opened with acrobat reader
%Please use one of the alternatives below to circumvent this error by uncommenting one or the other
%\pdfobjcompresslevel=0
%\pdfminorversion=4

% See the \addtolength command later in the file to balance the column lengths
% on the last page of the document

% The following packages can be found on http:\\www.ctan.org
% \usepackage{graphics} % for pdf, bitmapped graphics files
%\usepackage{epsfig} % for postscript graphics files
%\usepackage{mathptmx} % assumes new font selection scheme installed
%\usepackage{times} % assumes new font selection scheme installed
% \usepackage{amsmath} % assumes amsmath package installed
% \usepackage{amssymb}  % assumes amsmath package installed

%  Additional packages

\let\proof\relax
\let\endproof\relax
\usepackage[T1]{fontenc}
\usepackage{graphicx}
\usepackage{soul}
\usepackage{xcolor}
\usepackage{colortbl}  % Required for column coloring
\usepackage[english]{babel}
\usepackage{amsthm}
\usepackage{amssymb}
\newtheorem{assumption}{\textbf{Assumption}}
\newtheorem{theorem}{\textbf{Theorem}} % Theorem numbered within sections
\newtheorem{lemma}[theorem]{\textbf{Lemma}}
\newtheorem{corollary}[theorem]{\textbf{Corollary}}
\newtheorem{proposition}{\textbf{Proposition}}
\newtheorem{definition}{\textbf{Definition}}
\newtheorem{remark}{\textbf{Remark}}

\usepackage{amsmath,amssymb,amsfonts}
\usepackage{algorithmic}
\usepackage{graphicx}
\usepackage{textcomp}
\usepackage{colortbl}  % Required for column coloring
\usepackage{xcolor}
\usepackage[colorlinks=true, linkcolor=blue, citecolor=blue, urlcolor=blue, hypertexnames=false]{hyperref}


\usepackage{amsmath,amssymb,amsthm}
\usepackage[T1]{fontenc}
% \usepackage{todonotes}

\newcommand{\Acl}{A_{\mathrm{cl}}}




\renewcommand{\qedsymbol}{\hbox{$\blacksquare$}}
\newcommand{\squeezeup}{\vspace{-2.5mm}}

\title{\LARGE \bf
Optimal Modified Feedback Strategies in LQ Games under Control Imperfections
}


\author{Mahdis Rabbani$^{1}$, Navid Mojahed$^{1}$, and Shima Nazari$^{1}$% <-this % stops a space
\thanks{*This work was not supported by any organization.}% <-this % stops a space
\thanks{$^{1}$Mahdis Rabbani, Navid Mojahed, and Shima Nazari are with the Mechanical and Aerospace Engineering,
        University of California, Davis, 1 Shields Ave., Davis, 95616 CA, USA
        {\tt\small \{mrabbani, nmojahed, snazari\}@ucdavis.edu}}%
}


\begin{document}

\addtolength{\topmargin}{2mm}   % Increase top margin by 5mm
% \addtolength{\textheight}{-2mm} % Reduce text height by 10mm

\maketitle
\thispagestyle{empty}
\pagestyle{empty}

\begin{abstract}
Game-theoretic approaches and Nash equilibrium have been widely applied across various engineering domains. However, practical challenges such as disturbances, delays, and actuator limitations can hinder the precise execution of Nash equilibrium strategies. This work investigates the impact of such implementation imperfections on game trajectories and players' costs in the context of a two--player finite--horizon linear quadratic (LQ) nonzero-sum game. Specifically, we analyze how small deviations by one player, measured or estimated at each stage affect the state trajectory and the other player’s cost. To mitigate these effects, we construct a compensation law for the influenced player by augmenting the nominal game with the measurable deviation dynamics. The resulting policy is
shown to be optimal within a causal affine policy class, and,
for sufficiently small deviations, it locally outperforms the
uncompensated equilibrium-derived feedback. Rigorous analysis
and proofs are provided, and the effectiveness of the proposed
approach is demonstrated through a representative numerical
example.
% \todo[inline]{Before: affect the state and cost function of the other player. To mitigate these effects, we propose an adjusted control policy that optimally compensates for the deviations under the stated information structure and can, under certain conditions, exploit them to improve performance. Rigorous mathematical analysis and proofs are provided, and the effectiveness of the proposed method is demonstrated through a representative numerical example.}



\end{abstract}
% \begin{abstract}
% We study finite-horizon, discrete-time linear–quadratic (LQ) nonzero-sum games when one player’s implemented control deviates from its feedback Nash equilibrium (FNE) command due to execution effects such as actuator lag. First, we derive a trajectory- and cost-sensitivity expansion that quantifies how an opponent’s stagewise deviations propagate through the time-varying closed loop and impact the other player’s cost. Next, assuming the deviation is measured (or estimated) online, we synthesize a \emph{compensated feedback} (CF) policy via an augmented Riccati recursion that acts on both the state and the deviation signal. The CF law enjoys a nonnegative gap identity and is never worse—often strictly better—than the unmodified equilibrium-derived feedback to second order in the deviation magnitude; it reduces to the nominal FNE action when the deviation vanishes. We also provide a specialization for slowly varying lag dynamics that yields closed-form gain structure and recovers the FNE state gain. A two-cart spring–damper example with a lagged actuator demonstrates the performance recovery and provides intuition.
% \end{abstract}


\section{Introduction}
\label{sec:introduction}

Game theory is a powerful analytical framework for modeling rational decision-making in complex multi-agent systems, particularly when interactions are strategic and agents’ decisions influence each other~\cite{shamma2020game,basar1999dynamic}. Such interactions often involve conflicting objectives, motivating the development of solution concepts to characterize equilibrium behavior. Among these, the Nash equilibrium is a well-established notion, defined as a set of strategies from which no participant can unilaterally deviate to achieve a better outcome~\cite{nash1950equilibrium}.

Thanks to its solid theoretical foundation, Nash equilibrium has been widely applied in diverse engineering domains, including robotics~\cite{chiu2021encoding}, traffic management~\cite{Oszczypała2023Nash}, and sensor networks~\cite{Zheng2023game}, as well as in emerging areas such as energy management~\cite{YANG2020216nashQ,abdolmohammadi2024sizing}, autonomous vehicle navigation~\cite{nazari2024powertrain, bouzidi2025interaction, abtahi2025multi}, and human–robot interaction~\cite{soltanian2025peerawarecostestimationnonlinear}.

Despite its wide adoption, implementing Nash strategies in real-world engineering systems faces significant challenges due to uncertainties, disturbances, delays, and actuator imperfections~\cite{van-den=broek2023,AMATO2002507_guaranteeing,ma2024ϵ}. Extensions such as the $\varepsilon$-equilibrium relax optimality by allowing bounded deviations from Nash strategies~\cite{TANAKA1991413}, enabling numerical approximations~\cite{NORTMANN20231760,alvarez2009urban} and supporting further studies on sensitivity and robustness~\cite{van-den=broek2023,AMATO2002507_guaranteeing,Jimenez01072006,tan2005sensitivity,hernandez2024ϵ}.

Robustness of Nash equilibrium to uncertainty has been studied from various perspectives. For example, model uncertainties in differential games were addressed in~\cite{van-den=broek2023,AMATO2002507_guaranteeing,hernandez2024ϵ}, where~\cite{van-den=broek2023} designed a finite-horizon feedback Nash strategy for Linear Quadratic (LQ) games, \cite{AMATO2002507_guaranteeing} extended the scope to infinite-horizon settings, and \cite{hernandez2024ϵ} relaxed the LQ structure via weaker convexity assumptions. Other works addressed different uncertainty sources:~\cite{Jimenez01072006} proposed adaptive control to mitigate measurement noise, and~\cite{tan2005sensitivity} analyzed the effect of cost parameter variations in discrete-time games.

Nonetheless, the sensitivity of Nash equilibrium to player-specific control deviations, arising from implementation imperfections, remains largely unexplored. Such deviations, e.g., due to actuator errors or delays, are common in engineering systems and can directly degrade performance. Moreover, the discrete-time setting has received little attention in this context. This work analytically investigates how such deviations influence system trajectories and player costs.

To address these effects, we construct a compensation law for
the influenced player that uses the nominal equilibrium structure
together with measurable deviation information. The resulting
controller is not posed as a new Nash equilibrium of the modified
interaction; rather, it is a structured compensation policy around
the nominal game that mitigates performance loss caused by
execution-time deviations.
% \todo[inline]{Before: To address these effects, we propose a modified feedback policy that both mitigates performance loss and exploits available deviation information. }
While prior adaptive multi-agent methods exist~\cite{guerrero2021openloop,wang2022risk}, they do not explicitly target compensation for execution-time control deviations of a specific player, which is addressed in this work.

This paper investigates finite-horizon, discrete-time LQ
games in which one player’s control inputs deviate during execution.
We analyze the resulting impact on joint trajectories and on the
other player’s cost. The main contributions are summarized below:
\begin{itemize}
    \item A sensitivity analysis for discrete-time games, deriving closed-form expressions that characterize how one player’s deviations propagate through the system dynamics and influence the other player’s cost.
    \item
    A compensation design for the influenced player, obtained
  by augmenting the nominal FNE dynamics with measurable deviation dynamics. The resulting policy is shown to be uniquely optimal over a causal affine policy class that includes the uncompensated reference feedback as a special case.
  % \todo[inline]{Before: An optimal modified feedback control policy that compensates for such deviations and ensures performance that is never worse—and often strictly better—than the original feedback Nash equilibrium.}
    \item  A local performance comparison showing that, for sufficiently small deviations, the proposed compensated feedback is no worse than the uncompensated equilibrium-derived feedback to second order, together with a representative numerical example illustrating the resulting performance recovery.
    % \todo[inline]{Before: A complete theoretical analysis, supported by a representative numerical example, that highlights the robustness of Nash equilibrium under execution imperfections.}
\end{itemize}

The remainder of the paper is organized as follows.
Section~\ref{sec: problem statement} presents the problem formulation and background.
Section~\ref{sec:error-on-nash-equilibria} develops the sensitivity analysis.
Section~\ref{sec: control policy design} introduces the modified feedback control policy.
Section~\ref{sec:numerics} demonstrates its effectiveness through a numerical example.
Finally, Section~\ref{sec: conclusion} concludes the paper.


\noindent \textbf{Notation}. $\mathbb{R}$ and $\mathbb{Z}$ denote the sets of real and integer numbers, respectively. For a matrix $A$, $A^\top$ denotes its transpose; for a square matrix $B$, $B \succ 0$ ($B \succeq 0$) means $B$ is positive definite (positive semi-definite), and $B^{-1}$ denotes its inverse. For $i \in N_p=\{1,2\}$, the subscript $_{-i}$ refers to $N_{p} \setminus \{i\}$, i.e., all indices except $i$. The superscript $^\star$ denotes the exact optimal solution. A tilde, as in $\tilde P_{i,k}$, denotes auxiliary Riccati-type matrices used to distinguish the coupled finite-horizon game recursion from standard Riccati matrices.

\section{Preliminaries}
\label{sec: problem statement}
We focus on finite-horizon linear quadratic (LQ) games~\cite{Engwerda2009lq} with discrete-time dynamics
\begin{equation}
    \label{eq:DT game dynamics}
    x_{k+1} = A x_k + B_1 u_{1,k} + B_2 u_{2,k},
\end{equation}
for $k \in \mathbb{K} = \{0,1,\dots, N-1\} \subset \mathbb{Z}$, where:
\begin{itemize}
    \item \(x_k = \begin{bmatrix} x_{1,k}^\top & x_{2,k}^\top \end{bmatrix}^\top
\in \mathbb{R}^{n_1+n_2}\) is the joint state vector obtained by concatenating the state vectors $x_{i,k} \in \mathbb{R}^{n_i}$ of Player $i \in \{1,2\}$, where $n_i$ is the dimension of Player~$i$'s state.
    \item \(u_{1,k} \in \mathbb{R}^{m_1}\) and \(u_{2,k} \in \mathbb{R}^{m_2}\) are the control inputs of Player~1 and~2, respectively, where $m_i$ is the dimension of Player~$i$'s control input.
    \item \(A \in \mathbb{R}^{(n_1+n_2)\times (n_1+n_2)}\) is the joint state transition matrix.
    \item \(B_1 \in \mathbb{R}^{(n_1+n_2)\times m_1}\) and \(B_2 \in \mathbb{R}^{(n_1+n_2)\times m_2}\) are the input matrices for Player~1 and~2, respectively.
\end{itemize}

Suppose Player~$i \in \{1,2\}$ minimizes its cost $J_i$ over an $N$-step horizon by choosing control inputs $u_{i,k}$ at each stage $k \in \mathbb{K}$. The strategy of Player~$i$ is $U_i = \{u_{i,k}\}_{k=0}^{N-1}$, and $U_{-i}$ denotes its counterpart’s strategy. The quadratic cost function of Player~$i$ is
\begin{equation}
    \label{eq:LQ cost function}
    J_i = \sum_{k=0}^{N-1} \left( x_k^\top Q_i x_k + u_{i,k}^\top R_i u_{i,k} \right) + x_N^\top Q_{i,N} x_N,
\end{equation}
where $Q_i, Q_{i,N} \succeq 0$ are the stage and terminal weighting matrices, and $R_i \succ 0$ is the control weighting matrix.

We use Nash equilibrium to mean a strategy profile $U_i$, where no player can reduce its cost by a unilateral deviation, i.e. $ J_i(x_0, U_i^\star, U_{-i}^\star) \le J_i(x_0, U_i, U_{-i}^\star)$~\cite{nash1950equilibrium}.
\begin{definition}
\label{def: definition of Nash}
Consider the LQ game \eqref{eq:DT game dynamics}–\eqref{eq:LQ cost function}.  
A pair of feedback laws
\[
u_{i,k}=-K^\star_{i,k}x_k,\qquad i\in\{1,2\},\ k\in\mathbb{K},
\]
constitutes a linear Feedback Nash Equilibrium (FNE) if, for each $i$, $u_{i,k}$ minimizes $J_i$ against the opponent’s feedback $u_{-i,k}=-K^\star_{-i,k}x_k$ over the horizon.
\end{definition}

The FNE gains $\{K^\star_{i,k}\}$ are obtained by a backward Riccati recursion coupled across players \cite{salizzoni2025bridging}.  
Let $\tilde P_{i,N}=Q_{i,N}$ and define, for each $k\in\{0,\dots,N-1\}$ (computed backward),
\[
A^{-i}_{\mathrm{cl},k} \;:=\; A-\sum_{j\neq i} B_j K_{j,k}.
\]
Given $\{\tilde P_{i,k+1}\}$ and $\{K_{-i,k}\}$, the matrices $\tilde P_{i,k}$ satisfy
\begin{equation}
\label{eq:coupled_riccati-nominal}
\begin{aligned}
\tilde P_{i,k}
= Q_i &+ (A^{-i}_{\mathrm{cl},k})^\top \tilde P_{i,k+1} A^{-i}_{\mathrm{cl},k}\\
&- (A^{-i}_{\mathrm{cl},k})^\top \tilde P_{i,k+1} B_i
\,\tilde H_k\,^{-1}\!
B_i^\top \tilde P_{i,k+1} A^{-i}_{\mathrm{cl},k},
\end{aligned}
\end{equation}
where $\tilde H_k :=R_i + B_i^\top \tilde P_{i,k+1} B_i$. The corresponding FNE feedback gains are
\begin{equation}
\label{eq:Ki_formula-nominal}
K^\star_{i,k}
= \tilde H_k^{-1} B_i^\top \tilde P_{i,k+1} A^{-i}_{\mathrm{cl},k}.
\end{equation}

Under the standard assumptions $Q_i,Q_{i,N}\succeq 0$, $R_i\succ 0$ and nonsingularity of the gain system  in Definition~\ref{def: definition of Nash} at each $k$, the backward sweep $k=N-1\to 0$ yields a unique FNE \cite[Sec.~6]{Engwerda2009lq}.


% \begin{definition}
% \label{def: definition of Nash}
% A strategy profile $\{U_i^\star, U_{-i}^\star\}$ is an open-loop Nash equilibrium of the game \eqref{eq:LQ cost function} under the dynamics \eqref{eq:DT game dynamics} if
% \begin{equation*}
%     J_i(x(0), U_i^\star, U_{-i}^\star) \le J_i(x(0), U_i, U_{-i}^\star)
% \end{equation*}
% for all admissible $U_i$, $i \in \{1,2\}$~\cite{nash1950equilibrium}.
% \end{definition}

We next recall standard sufficient conditions that ensure the existence and uniqueness of the Nash equilibrium in Definition~\ref{def: definition of Nash}, thereby guaranteeing well-posedness of the game.

\begin{remark}
Consider the following sufficient conditions for the existence of a unique Nash equilibrium in a two-player finite-horizon nonzero-sum discrete-time LQ game \eqref{eq:LQ cost function}:
\begin{enumerate}
    \item The state weighting matrices satisfy $Q_{i}, Q_{i,N} \succeq 0$ and the control weighting matrices satisfy $R_{i} \succ 0$ for $i \in \{1,2\}$.\label{condition 1}
    \item Let $\tilde{B} = \begin{bmatrix} B_1 & B_2 \end{bmatrix}$. The pair $(A, \tilde{B})$ is controllable.\label{condition 2}
    % , i.e., the controllability matrix
    % \[
    % \mathcal{C} = \begin{bmatrix} \bar{B} & A\bar{B} & \dots & A^{n-1} \bar{B} \end{bmatrix},
    % \]
    % with $n = n_1 + n_2$, has full rank.
    \item Define
    \[
    M_{k} =
    \begin{bmatrix}
    I & F_1 & F_2 \\
    P_{1,k} & -I & 0 \\
    P_{2,k} & 0 & -I
    \end{bmatrix},
    \quad F_i = B_i R_{i}^{-1} B_i^\top,
    \]
    where $I$ is the $n\times n$ identity and $P_{i,k}$ denote the player Riccati matrices from the standard backward recursions; invertibility of $M_k$ ensures well-posed coupled updates.\label{condition 3}
\end{enumerate}
Under conditions~\ref{condition 1}-\ref{condition 3}, a unique finite-horizon Nash equilibrium exists. We assume these conditions henceforth \cite{LAA_Jungers_2013}.
\end{remark}


\section{First-Order Sensitivity of Player~1 Cost under Opponent Deviations}
\label{sec:error-on-nash-equilibria}
Under the memoryless state-feedback information structure, a finite-horizon
FNE for a two-player game has the form
\begin{equation}
\label{eq:feedback nash definition}
u^\star_{i,k} = -K^\star_{i,k} x_k, \qquad i\in\{1,2\},\ \ k=0,\dots,N-1,
\end{equation}
as stated in Definition~\ref{def: definition of Nash}.
% with gains $\{K^\star_{1,k},K^\star_{2,k}\}$ arising from coupled Riccati
% recursions; see, e.g.,~\cite{nortmann2023approximate}. Substituting
Substituting \eqref{eq:feedback nash definition} into the dynamics \eqref{eq:DT game dynamics} yields the
time-varying closed-loop matrix
\begin{equation}
\label{eq:Acl-def}
A_{\mathrm{cl},k} := A - B_1 K^\star_{1,k} - B_2 K^\star_{2,k},
\end{equation}
which evolves the dynamics
\begin{equation}
\label{eq:closed-loop Nash dynamic}
x_{k+1}^\star = A_{\mathrm{cl},k} x_k^\star,\qquad x_0^\star \equiv x_0.
\end{equation}

In practice, Player~2’s implemented input may deviate from its FNE command,
$u_{2,k} = u^\star_{2,k} + \Delta u_{2,k}$. In this section, we quantify the
effect of small deviations $\{\Delta u_{2,k}\}$ on Player~1’s cost when
Player~1 keeps its FNE law \eqref{eq:feedback nash definition}. We obtain
(i) an exact expression for the induced state perturbation and
(ii) a first-order expansion of $\Delta J_1$ with explicit coefficients that
respect the time variation in \eqref{eq:Acl-def}.


%-----------------------------
\begin{definition}
\label{def:phi}
For integers $l\ge k\ge j$, define
\begin{equation*}
\Phi(k,j) := \prod_{i=j}^{k-1} A_{\text{cl},i} \;=\; A_{\text{cl},k-1} \cdots A_{\text{cl},j},
\end{equation*}
with the convention $\Phi(j,j):=I$. 
% Then,
% \[
% \Phi(k+1,j)=A_{\text{cl},k}\,\Phi(k,j), \qquad\Phi(k,\ell)\,\Phi(\ell,j)=\Phi(k,j).
% \]
\end{definition}
The state-transition matrix satisfies the standard identities
\[
\Phi(k+1,j)=A_{\text{cl},k}\,\Phi(k,j), \qquad\Phi(k,\ell)\,\Phi(\ell,j)=\Phi(k,j).
\]
% \todo[inline]{Reviewer 2, Comment 2}

%-----------------------------
% Proposition 1
%-----------------------------
\begin{proposition}\label{prop:first_order}
Let $\{K^\star_{1,k}, K^\star_{2,k}\}_{k=0}^{N-1}$ be the FNE gains in \eqref{eq:feedback nash definition}. Suppose Player~2 deviates as $u_{2,k}=u^\star_{2,k}+\Delta u_{2,k}$ while Player~1 keeps $u_{1,k}=u^\star_{1,k}=-K^\star_{1,k}x_k$. Then, for $k=1,\dots,N$,
\begin{equation}\label{eq:dx_formula}
\Delta x_k \;=\; \sum_{j=0}^{k-1} \Phi(k,j+1)\,B_2\,\Delta u_{2,j},
\end{equation}
and, with $S_k := Q_1 + K_{1,k}^{\star\top} R_1 K^\star_{1,k}$,
\begin{equation}\label{eq:DJ_linear}
\Delta J_1 \;=\; 2 \sum_{j=0}^{N-1} x_0^\top \Lambda_j \,B_2\,\Delta u_{2,j} \;+\; \mathcal{O}\!\big(\|\Delta u_2\|_2^2\big),
\end{equation}
where
\begin{equation}\label{eq:Lambda_j}
\begin{aligned}
    \Lambda_j \;=\; \sum_{k=j+1}^{N-1} \Phi(k,0)^\top &S_k \,\Phi(k,j+1) \\
    \;&+\; \Phi(N,0)^\top Q_{1,N}\,\Phi(N,j+1),
\end{aligned}
\end{equation}
for $j=0,\dots,N-1$.
Here $\Phi(k,j)$ is as in Definition~\ref{def:phi}, $\Delta u_2:=[\Delta u_{2,0}^\top\,\cdots\,\Delta u_{2,N-1}^\top]^\top$, and
the remainder is with respect to the Euclidean norm $\|\Delta u_2\|_2$.
\end{proposition}

\begin{proof}
Under the FNE
% and defining $x_0^\star \equiv x_0$
, the nominal trajectory is obtained via \eqref{eq:closed-loop Nash dynamic}. With $u_{2,k}=u_{2,k}^\star+\Delta u_{2,k}$ and $u_{1,k}=u_{1,k}^\star$, the actual state satisfies
\begin{equation*}
x_{k+1} \;=\; A_{\text{cl},k} \, x_k + B_2 \, \Delta u_{2,k}.
\end{equation*}
Defining $\Delta x_k := x_k - x_k^\star$ gives
\begin{equation*}\label{eq:dx_recursion}
\Delta x_{k+1} \;=\; A_{\text{cl},k} \, \Delta x_k + B_2 \, \Delta u_{2,k}, \qquad \Delta x_0 = 0,
\end{equation*}
whose iteration yields \eqref{eq:dx_formula} using $\Phi(k+1,j)=A_{\text{cl},k}\,\Phi(k,j)$.

Next, linearizing \eqref{eq:LQ cost function} for $i=1$ around the FNE trajectory $(x_k^\star,u_{1,k}^\star)$, leads to
\begin{equation*}\label{eq:DJ_raw}
\begin{aligned}
    \Delta J_1 \;=\; 2 \sum_{k=0}^{N-1} \big( x_k^{\star\top}& Q_1 \Delta x_k + u_{1,k}^{\star\top} R_1 \Delta u_{1,k} \big)\\
    \;+\; &2 \, x_N^{\star\top} Q_{1,N} \Delta x_N \;+\; \mathcal{O}(\|\Delta x\|_2^2),
\end{aligned}
\end{equation*}
where $\Delta\,J_1 = J_1(x_k,u_{1,k})-J_1(x_k^\star,u_{1,k}^\star)$.

Since $u_{1,k}=-K^\star_{1,k}x_k$, we have $\Delta u_{1,k} = -K^\star_{1,k}\Delta x_k$, hence
\begin{equation}\label{eq:DJ_Sk}
\begin{aligned}
    \Delta J_1 \;=\; 2 \sum_{k=0}^{N-1} x_k^{\star\top} S_k \, \Delta x_k
\;+\; 2 x_N^{\star\top}& Q_{1,N} \Delta x_N \\
\;&+\; \mathcal{O}(\|\Delta x\|_2^2),
\end{aligned}
\end{equation}
where $S_k := Q_1 + K_{1,k}^{\star\top} R_1 K^\star_{1,k}$. Substituting \eqref{eq:dx_formula} into \eqref{eq:DJ_Sk}, defining $\Lambda_j$ as in \eqref{eq:Lambda_j}, and using $x_k^\star=\Phi(k,0)x_0$ gives \eqref{eq:Lambda_j}. As $\Delta x$ is linear in $\Delta u_2$ by \eqref{eq:dx_formula}, the remainder is $\mathcal{O}(\|\Delta u_2\|_2^2)$. This completes the proof.
\end{proof}

% \begin{theorem}
% \label{thm:CostAndStateDev}
% \hl{change to proposition}

% Let $\{U_1^\star,U_2^\star\}$ be the feedback Nash equilibrium of the game \eqref{eq: LQ cost function} with the system dynamics \eqref{eq: DT game dynamics} and the feedback Nash policy \eqref{eq: feedback nash definition}. Suppose that Player~2 deviates from its feedback Nash strategy according to \eqref{eq: u2 with uncertainty}. Then, for small control deviations, the effects of $\Delta u_{2,k}$ propagate through the state trajectory and Player~1's cost function according to the following linear mappings:
% \begin{equation}
%     \label{eq: error propagation in state}
%     \Delta x_k 
%     = \sum_{j=0}^{k-1} 
%     \left( \prod_{t=j+1}^{k-1} A_{\mathrm{cl},t} \right) B_2 \,\Delta u_{2,j},
% \end{equation}
% and
% \begin{align}
%     \label{eq: error propagation in p1's cost}
%     \Delta J_1 \approx 2\,x_0^\top \sum_{j=0}^{N-1} \lambda_j\,\Delta u_{2,j},
% \end{align}
% where
% \[
% \lambda_j = \sum_{k=j+1}^{N} 
% \left( \prod_{t=0}^{k-1} A_{\mathrm{cl},t} \right)^\top B_2
% \, S_k \left( \prod_{t=j+1}^{k-1} A_{\mathrm{cl},t} \right),
% \]
% \[
% S_k = Q_1 + K_{1,k}^{\star\top} R_1 K_{1,k}^\star.
% \]


% and $A_{\mathrm{cl},k}$ is the time-varying closed-loop matrix defined in \eqref{eq: closed-loop Nash dynamic}.
% \end{theorem}

% \begin{proof}
% Substituting Player~1's feedback strategy from \eqref{eq: feedback nash definition} 
% and Player~2's perturbed strategy \eqref{eq: u2 with uncertainty} into the system 
% dynamics \eqref{eq: DT game dynamics} yields the perturbed closed-loop dynamics:
% \begin{equation}
% \label{eq: perturbed state dynamics}
%     x_{k+1} = A_{\mathrm{cl},k} \, x_k + B_2\,\Delta u_{2,k},
% \end{equation}
% where \(A_{\mathrm{cl},k} := A - B_1 K_{1,k}^\star - B_2 K_{2,k}^\star\).

% Let the state deviation be \(\Delta x_k := x_k - x_k^\star\).  
% Subtracting \eqref{eq: closed-loop Nash dynamic} from 
% \eqref{eq: perturbed state dynamics} gives the error dynamics
% \begin{equation}
%     \label{eq: error dynamics}
%     \Delta x_{k+1} = A_{\mathrm{cl},k}\,\Delta x_k + B_2\,\Delta u_{2,k}.
% \end{equation}
% Under the assumption \(\Delta x_0 = 0\), backward expansion of 
% \eqref{eq: error dynamics} yields:
% \begin{equation}
% \label{eq: error propagation in state}
%     \Delta x_k = \sum_{j=0}^{k-1} 
%       \left( \prod_{t=j+1}^{k-1} A_{\mathrm{cl},t} \right) B_2\,\Delta u_{2,j},
% \end{equation}
% where the empty product for \(j = k-1\) is taken as the identity.

% ---

% Next, consider the deviation in Player~1’s cost function.  
% Let \(\Delta J_1 := J_1 - J_1^\star\).  
% Substituting \(x_k = x_k^\star + \Delta x_k\) into \eqref{eq: LQ cost function} 
% for \(i=1\) and discarding second-order terms in the deviations gives:
% \begin{equation}
% \label{eq: cost linearized}
% \begin{aligned}
% \Delta J_1
%     &\approx
% \sum_{k=0}^{N-1} 
% \left[
%   2\,x_k^{\star\top} Q_1\,\Delta x_k
%  +2\,u_{1,k}^{\star\top} R_1\,\Delta u_{1,k}
% \right] \\
% &\quad + 2\,x_N^{\star\top} Q_{1,N}\,\Delta x_N.
% \end{aligned}
% \end{equation}

% From \(u_{1,k}^\star = -K_{1,k}^\star\,x_k^\star\) and 
% \(\Delta u_{1,k} = -K_{1,k}^\star\,\Delta x_k\), we have:
% \[
% u_{1,k}^{\star\top} R_1\,\Delta u_{1,k}
%     = x_k^{\star\top} K_{1,k}^{\star\top} R_1 K_{1,k}^\star\,\Delta x_k.
% \]
% Defining
% \[
% S_k := Q_1 + K_{1,k}^{\star\top} R_1 K_{1,k}^\star,
% \]
% equation \eqref{eq: cost linearized} becomes:
% \[
% \Delta J_1
% \approx
% 2\sum_{k=0}^{N-1} x_k^{\star\top} S_k\,\Delta x_k
%  + 2\,x_N^{\star\top} Q_{1,N}\,\Delta x_N.
% \]

% Substituting \eqref{eq: error propagation in state} into this expression and 
% interchanging the order of summations, we get:
% \begin{align*}
% \Delta J_1
% &\approx
% 2\sum_{j=0}^{N-1}
% \left(
%     \sum_{k=j+1}^{N-1}
%       x_k^{\star\top} S_k 
%       \left( \prod_{t=j+1}^{k-1} A_{\mathrm{cl},t} \right) B_2
% \right) \Delta u_{2,j} \\
% &\quad + 2\,x_N^{\star\top} Q_{1,N}\,\Delta x_N.
% \end{align*}

% Using \(x_k^\star = \left( \prod_{t=0}^{k-1} A_{\mathrm{cl},t} \right) x_0\), 
% the inner product terms can be expressed in closed form, which yields 
% \eqref{eq: error propagation in p1's cost}.  
% This completes the proof.
% \end{proof}

% =====================================================
% Stability Remark
% =====================================================
% \begin{remark}
%     Suppose $A_{\mathrm{cl}}$ is stable, meaning that its eigenvalues lie within the unit disk, and the control deviation is persistent but bounded. In that case, the system satisfies Bounded-Input Bounded-Output (BIBO) stability, ensuring that the state deviation remains bounded for all $k$. However, unless $\Delta u_{2,k} \to 0$ as $k \to \infty$, the state deviation does not necessarily decay to zero, meaning that the error dynamics are not asymptotically stable.
% \end{remark}

% \begin{remark}
%     Theorem~\ref{thm:CostAndStateDev} provides insights into sensitivity of feedback Nash strategies to small control deviations arising from implementation imperfections. Representative application of such imperfections is in robotic systems, etc.
%     It is noteworthy that Theorem~\ref{thm:CostAndStateDev} assumes that despite the perturbation, both players stick to their feedback Nash policy and use the same feedback gains as in the non-perturbed case, meaning that in this theorem, there is no adjustment neither in the policy, nor in the gains. This invariance of \( K \) is crucial to maintaining the closed-loop stability of the system, as it ensures that the core feedback strategy derived from the Nash equilibrium remains valid. 
% \end{remark}

% \hl{delete the next remark or clarify it.}
% \begin{remark}
%     The cumulative formulation of $\lambda_j$ as the weighting matrix of the deviation indicates that the perturbation keeps influencing the cost function deviation over all subsequent time steps. However, this impacts decays over time according to the closed-loop dynamics defined by $A_{\mathrm{cl}}$.
% \end{remark}

% As a result of the linear propagation stated in Theorem~\ref{thm:CostAndStateDev}, Player~1 can compensate for the error by adjusting its Nash strategy policy as stated in Theorem~\ref{prop:OptimalActionP1}. 

Proposition~\ref{prop:first_order} shows the deviations of Player~2 linearly perturb Player~1’s cost. If Player~1 keeps the equilibrium-derived feedback without compensation, e.g. $u_{1,k}^{\mathrm{REF}}:=-K^\star_{1,k}x_k$ (referred to as \textit{Reference Feedback (REF)}), its policy is generally sub-optimal under the realized dynamics. Motivated by this sensitivity, we next design a controller that models the deviation and attenuates its effect as it propagates through the closed loop, compensating for these deviations.


\section{Control Policy Design with Disturbance Compensation}
\label{sec: control policy design}

In practice, the control implementation is rarely exact. Differences between the commanded and implemented inputs typically arise from (i) actuator/servo dynamics and inner-loop tracking (first-order lags, rate/slew limits), (ii) computation/communication latency and sample-and-hold effects, and (iii) command filtering/smoothing or safety layers that reshape inputs.
These effects induce an input disturbance $w_k \;:=\; B_2\big(u_{2,k}-u_{2,k}^\star\big)$, acting precisely through the channel identified by the sensitivity result.

In practice, $w_k$ need not be directly measured at the actuator. Instead, Player~1 may form an estimate $\hat w_k$ from available information. If the commanded input $u^\star_{2,k}=-K^\star_{2,k}x_k$ is known from the nominal FNE model and the implemented input $u_{2,k}$ can be communicated, sensed, or reconstructed from actuator or vehicle telemetry, then one may compute $\hat w_k = B_2\bigl(u_{2,k}-u^\star_{2,k}\bigr)$.
When $u_{2,k}$ is not directly available, $\hat w_k$ can instead be obtained using a standard disturbance observer or state estimator built from the plant model and measured state trajectory. In that case, the proposed CF law is applied using $\hat w_k$ in place of $w_k$, while the present analysis corresponds to the ideal full-information case.
% \todo[inline]{Reviewer 1, Comment 2.}

We therefore model $w_k$ using simple, identifiable dynamics and equip Player~1 with a \textit{Compensated Feedback (CF)} law containing a feedforward term $L_k w_k$ to offset the predictable component of the disturbance. In this way, the feedback term $K_k$ regulates the residual error rather than reacting only after the state has already deviated. We first consider a physically grounded actuator-lag model and then narrow it down to a slow-varying error model.
% \todo[inline]{Before: We, therefore, model $w_k$ with simple, identifiable dynamics. Then, we equip Player~1 with a \textit{Compensated Feedback (CF)} law, which includes a feedforward term $L_kw_k$ to cancel the predictable component of the disturbance at the actuator, so the feedback term $K_k$ regulates the residuals instead of reacting after the state has already moved. We first treat a physically grounded actuator-lag model, then narrow it down to a specific case of slow-varying error dynamics. }


\subsection{Compensated Feedback under Lagged Execution Error Dynamics}
\label{sec:lag-best-response}

Assume that Player~2’s actuator is a discrete-time first-order lag that tracks its desired (Nash) input $u^{\star}_{2,k}=-K^{\star}_{2,k}x_k$ and let $u_{2,k}$ denote the implemented command. Let the execution error be $e_k:=u_{2,k}-u^{\star}_{2,k}$. Then, the plant disturbance is $w_k=B_2 e_k$.

\begin{assumption}
\label{ass:info-lag}
At time $k$, Player~1 observes $(x_k,w_k)$ and Player~2 implements its FNE command $u^{\star}_{2,k}=-K^{\star}_{2,k}x_k$ through a discrete first-order lag:
\begin{equation}\label{eq:u2-lag}
u_{2,k+1} \;=\; \alpha\,u_{2,k} \;+\; (1-\alpha)\,u^{\star}_{2,k}, \qquad \alpha\in(0,1).
\end{equation}
The induced disturbance recursion is
\begin{equation}\label{eq:w-lag}
\begin{aligned}
w_{k+1} &\;=\; \alpha\,w_k \;-\; B_2\big(u^{\star}_{2,k+1}-u^{\star}_{2,k}\big)\\
&\;=\; \alpha\,w_k \;+\; B_2\!\big(K^{\star}_{2,k+1}x_{k+1}-K^{\star}_{2,k}x_k\big).
\end{aligned}
\end{equation}
It is noteworthy that at each stage $k$, Player 1 knows the current $x_k$ and the current $w_k$, but not future values.
\end{assumption}

\begin{remark}
\label{rem:alpha-tau}
The discrete recursion \eqref{eq:u2-lag} is the zero-order-hold (ZOH) sampled counterpart of the continuous-time actuator $\dot u_2(t)=\tfrac{1}{\tau}(u_2^{\star}(t)-u_2(t))$ with sampling period $\Delta t$. The pole maps as
\begin{equation}\label{eq:alpha-tau}
\alpha \;=\; e^{-\Delta t/\tau},
\end{equation}
so the time constant $\tau$ of the continuous model sets the per-sample settling factor $\alpha\in(0,1)$.
\end{remark}


The lag model stated in Assumption~\ref{ass:info-lag}  captures two ubiquitous effects: (i) if the target $u^{\star}_{2,k}$ is steady, the tracking error decays geometrically with factor $\alpha$; (ii) changes in $u^{\star}_{2,k}$ inject a transient error that then dies out. Such behavior is widely used as a realistic low-order approximation of implementation effects.


% ******************************************************
Fix Player~2 at its Nash feedback $u^{\star}_{2,k}=-K^{\star}_{2,k}x_k$ and define
\begin{equation}\label{eq:Aeff-def}
A_{\mathrm{eff},k}:=A-B_2K^{\star}_{2,k},\qquad
x_{k+1}=A_{\mathrm{eff},k}x_k+B_1u_{1,k}+w_k .
\end{equation}

\begin{theorem}
\label{thm:lag-optimal}
Consider the game \eqref{eq:LQ cost function} with dynamics \eqref{eq:DT game dynamics}.
Let Player~2 obey its FNE policy \eqref{eq:feedback nash definition} and define the effective matrices
$A_{\mathrm{eff},k}$ via \eqref{eq:Aeff-def}. Under Assumption~\ref{ass:info-lag}, let the disturbance $w_k$ be generated by the lag recursion
\eqref{eq:w-lag}. Assume $Q_1,Q_{1,N}\succeq 0$ and $R_1\succ 0$. Over the causal affine class
\begin{equation}\label{eq:policy-class}
u_{1,k} \;=\; -K_k x_k \;-\; L_k w_k , \qquad k=0,\dots,N-1,
\end{equation}
there exists a \emph{unique} optimal policy that minimizes \eqref{eq:LQ cost function}, with gains
\begin{equation}\label{eq:gains-closed}
\big[\,K_k\ \ L_k\,\big]
\;=\;
\big(R_1+\bar B_k^\top \bar P_{k+1}\bar B_k\big)^{-1}\,\bar B_k^\top \bar P_{k+1}\bar A_k,
\end{equation}
where $\{\bar P_k\}$ is the unique sequence generated by the finite-horizon Riccati recursion
\begin{equation}\label{eq:riccati-aug}
\begin{aligned}
&\bar P_N=\bar Q_N,\\
&\begin{aligned}
\bar P_k&=\bar Q+\bar A_k^\top \bar P_{k+1}\bar A_k\\
&-\bar A_k^\top \bar P_{k+1}\bar B_k\big(R_1+\bar B_k^\top \bar P_{k+1}\bar B_k\big)^{-1}\bar B_k^\top \bar P_{k+1}\bar A_k,
\end{aligned}\\
&\text{for } k=N-1,\dots,0,
\end{aligned}
\end{equation}
on the augmented system with matrices
\begin{equation}\label{eq:ABQ-def}
\begin{aligned}
&\bar A_k:=\begin{bmatrix}
A_{\mathrm{eff},k} & I\\[2pt]
B_2K^{\star}_{2,k+1}A_{\mathrm{eff},k}-B_2K^{\star}_{2,k} & \alpha I + B_2K^{\star}_{2,k+1}
\end{bmatrix},\\
&\bar B_k:=\begin{bmatrix}B_1\\[2pt] B_2K^{\star}_{2,k+1}B_1\end{bmatrix},
\end{aligned}
\end{equation}
where
\begin{equation}\label{eq:Q-aug}
    \bar{Q} = \begin{bmatrix} Q_1 & 0 \\ 0 & 0 \end{bmatrix}, \quad \bar{Q}_N = \begin{bmatrix} Q_{1,N} & 0 \\ 0 & 0 \end{bmatrix}.
\end{equation}
Moreover, for any causal $\{\tilde K_k,\tilde L_k\}$, the incurred cost satisfies the nonnegative gap identity
\begin{equation}\label{eq:gap}
\begin{aligned}
\sum_{k=0}^{N-1}\!\big(\tilde u_{1,k}-u_{1,k}^\star\big)&^\top
\big(R_1+\bar B_k^\top \bar P_{k+1}\bar B_k\big)
\big(\tilde u_{1,k}-u_{1,k}^\star\big)\\
&\,=\,J_1(\{\tilde K_k,\tilde L_k\})-J_1(\{K_k,L_k\})\ \ge 0,
\end{aligned}
\end{equation}
where $u_{1,k}^\star:=-K_k x_k - L_k w_k$ and $\tilde u_{1,k}:=-\tilde K_k x_k - \tilde L_k w_k$.
\end{theorem}


\begin{proof}
Inspired by the Disturbance-Based State Compensation (DBSC) approach \cite{rugh1996linear}, which is commonly used in traditional state-feedback control design, and given Assumption~\ref{ass:info-lag}, we define the augmented state as $z_k:=[x_k^\top\ w_k^\top]^\top$.

Then the game problem can be rewritten in augmented form
\begin{align}
&\min_{ \{ u_{1,k} \}_{k=0}^{N-1} } z_N^\top \bar Q_N z_N + \sum_{k=0}^{N-1}\!\big(z_k^\top \bar Q z_k + u_{1,k}^\top R_1 u_{1,k}\big), \label{eq:aug-cost}\\
& \text{subject to:} \quad z_{k+1}=\bar A_k z_k + \bar B_k u_{1,k}, \quad k =0,\dots, N-1,\label{eq:augAB}
\end{align}
with $\bar A_k,\bar B_k,\bar Q,\bar Q_N$ given in \eqref{eq:ABQ-def}.

Define the value function of the augmented system
\begin{equation}\label{eq:Vk}
V_k(z):=\min_{\{u_{1,\ell}\}_{\ell=k}^{N-1}}
\Big\{z_N^\top \bar Q_N z_N+\sum_{\ell=k}^{N-1}\!(z_\ell^\top \bar Q z_\ell + u_{1,\ell}^\top R_1 u_{1,\ell})\Big\}
\end{equation}
subject to the augmented dynamic constraint of \eqref{eq:augAB}.

Based on the Bellman principle of optimality at time $k$, we can reformulate $V_k(z)$ as
\begin{equation}
\label{eq:V_k vs V_k+1}
    V_k(z_k)=\min_{u_{1,k}}\Big\{ z_k^\top \bar{Q} z_k + u_{1,k}^\top R_1 u_{1,k} + V_{k+1}\big(z_{k+1}\big) \Big\},
\end{equation}

% Backward induction shows $V_k(z)=z_k^\top \bar P_k z_k + \bar c_k$, with base $\bar P_N=\bar Q_N$. 
Backward induction yields $V_k(z_k)=z_k^\top \bar P_k z_k + \bar c_k$ with $\bar P_k$ given by \eqref{eq:riccati-aug}.
Substituting the augmented dynamics \eqref{eq:augAB} in \eqref{eq:V_k vs V_k+1} and collecting terms in $u_{1,k}$ yields the strictly convex quadratic in $u_{1,k}$:
\begin{equation}
\begin{aligned}
\mathcal Q_k(z_k,u_{1,k})
= u_{1,k}^\top &\mathcal H_k\, u_{1,k}
+ 2\,u_{1,k}^\top \bar B_k^\top \bar P_{k+1}\bar A_k\, z_k\\
&+ z_k^\top\!\big(\bar Q+\bar A_k^\top \bar P_{k+1}\bar A_k\big) z_k
+ \bar c_{k+1},
\end{aligned}
\end{equation}
where $\mathcal H_k:=R_1+\bar B_k^\top \bar P_{k+1}\bar B_k$ which is positive definite as $R_1\succ 0$ and $\bar P_{k+1}\succeq 0$.


Hence, the unique minimizer is
% \begin{equation}\label{eq:ukstar}
% u_{1,k}^\star \;=\; -\big(\underbrace{R_1+\bar B_k^\top \bar P_{k+1}\bar B_k}_{\mathcal H_k}\big)^{-1}\bar B_k^\top \bar P_{k+1}\bar A_k\, z_k,
% \end{equation}
\begin{equation}\label{eq:ukstar}
u_{1,k} \;=\; -\mathcal H_k^{-1}\bar B_k^\top \bar P_{k+1}\bar A_k\, z_k =:-\big[\,K_k\ \ L_k\,\big] z_k
\end{equation}
which, upon partitioning $z_k=[x_k^\top \; w_k^\top]^\top$, gives \eqref{eq:gains-closed}. 
% which results in the affine gains \eqref{eq:gains-closed} after partitioning based on $z_k=[x_k^\top\ w_k^\top]^\top$.

Completing the square yields \eqref{eq:riccati-aug}. Summing the stagewise completion terms from $k=0$ to $N-1$ telescopes the value functions and proves \eqref{eq:gap}. Uniqueness follows the strict convexity ($\mathcal H_k\succ 0$).
\end{proof}

% =======================================================
% Terminology for the FNE and REF and CF
% =======================================================
% \paragraph*{Terminology.}
% We denote by \emph{Feedback Nash Equilibrium (FNE)} the nominal equilibrium pair
% $\{K^{\star}_{1,k},K^{\star}_{2,k}\}_{k=0}^{N-1}$.
% As a baseline for Player~1 we use the \emph{Reference Feedback (REF)} controller
% \[
% u_{1,k}^{\mathrm{REF}} := -K^{\star}_{1,k} x_k,
% \]
% which applies the equilibrium-derived gain and does not use the measured disturbance.
% Our proposed \emph{Compensated Feedback (CF)} controller augments state feedback with an
% optimal preview/feedforward term:
% \[
% u_{1,k}^{\mathrm{CF}} := -K_k x_k - L_k w_k .
% \]
% In the results below, CF is the unique optimal best response for Player~1 under the assumed
% information structure.

% Intuition: the feedforward path cancels (or attenuates) the known effect of the opponent’s deviation before it propagates through the plant, while the feedback path cleans up the residual state error. This is the standard preview/disturbance-accommodation principle in optimal control.


By \eqref{eq:gap} in Theorem~\ref{thm:lag-optimal}, the gains $\{K_k, L_k\}$ are optimal over the causal affine class \eqref{eq:policy-class}, which includes the REF baseline $u_{1,k}^{\text{REF}}$ as the special case $L_k\equiv0$. 
Consequently, when execution-time deviations by Player 2 are measurable, Player~1 can improve upon the uncompensated reference feedback by adopting the compensated law within this policy class. This does not redefine the Nash equilibrium of the underlying game: the FNE is associated with the nominal game without implementation deviations, whereas the proposed law is a compensation strategy for the realized disturbed interaction.
% \todo[inline]{Before: Consequently, execution-time deviations by Player~2 create an opportunity for Player~1 to outperform its own equilibrium-derived feedback (REF) by adopting the compensated law. This does not contradict the definition of a Nash equilibrium: the FNE is defined for the nominal game without implementation deviations, whereas here Player~2’s realized input differs from its Nash command.}
The approach is practically relevant because such deviations routinely arise from disturbances, physical limits, and execution imperfections.

The next result quantifies the local benefit, giving a signed, second-order improvement in the deviation magnitude.

\begin{corollary}
\label{cor:local-gain-lag}
Let $J_1^{\mathrm{CF}}$ and $J_1^{\mathrm{REF}}$ be the costs under
$u_{1,k}^{\mathrm{CF}}=-K_k x_k - L_k w_k$ and
$u_{1,k}^{\mathrm{REF}}=-K^{\star}_{1,k} x_k$, respectively, where
$\{\bar P_k\}$ and $[\,K_k\ \ L_k\,]$ are as in Theorem~\ref{thm:lag-optimal}.
Define the selector $E_w:=[\,0\ \ I\,]$ so that $w_k=E_w z_k$ for
$z_k:=[x_k^\top\ w_k^\top]^\top$. For sufficiently small deviations
$\{\Delta u_{2,k}\}$ where $w_k=\mathcal{O}(\|\Delta u_2\|_2)$,
\begin{equation}\label{eq:local-gain-lag}
\begin{aligned}
J_1^{\mathrm{CF}} - J_1^{\mathrm{REF}}
\;=\;
-\,\sum_{k=0}^{N-1} \Big\| \mathcal{H}_k^{-\tfrac12}
\,\bar B_k^\top &\bar P_{k+1} E_w\, w_k \Big\|_2^{\,2}\\
&\;+\; \mathcal{O}\!\big(\|\Delta u_2\|_2^3\big)
\;\le\; 0,
\end{aligned}
\end{equation}
where $\displaystyle \mathcal{H}_k:=R_1+\bar B_k^\top \bar P_{k+1}\bar B_k \succ 0$. Thus, CF never performs worse than REF to second order in the deviation magnitude and strictly improves unless $w_k\equiv0$.
\end{corollary}
\begin{proof}
By Theorem~\ref{thm:lag-optimal}, the stagewise completion of squares gives, for any
candidate $\tilde u_{1,k}$,
$\mathcal{Q}_k(z_k,\tilde u_{1,k})-\mathcal{Q}_k(z_k,u_{1,k}^\star)
=\|\mathcal{H}_k^{1/2}(\tilde u_{1,k}-u_{1,k}^\star)\|_2^2$.
Evaluating this at the REF choice $\tilde u_{1,k}=-K_{1,k}^\star x_k$ yields
$\tilde u_{1,k}-u_{1,k}^\star= L_k w_k$ plus higher-order state terms induced by
$\Delta u_2$, which are $\mathcal{O}(\|\Delta u_2\|_2^2)$. Summing over $k$ gives
\eqref{eq:local-gain-lag} with the remainder $\mathcal{O}(\|\Delta u_2\|_2^3)$ and completes this proof.
\end{proof}



% \begin{proof}
% Start from the exact stagewise gap identity (completion of squares) comparing CF with an arbitrary competitor, and specialize the competitor to the REF baseline. Write
% $u_{1,k}^{\mathrm{REF}}-u_{1,k}^{\mathrm{CF}}=(K_k-K^{\star}_{1,k})x_k + L_k w_k$,
% substitute the Bellman-optimal gains, and use the sensitivity result that
% $w_k=\mathcal{O}(\Delta u_2)$ while the residual state terms are $\mathcal{O}(\Delta u_2)$.
% The leading nonzero contribution is the negative quadratic in \eqref{eq:local-gain};
% the remainder is cubic in $\|\Delta u_2\|$.
% \end{proof}

% \hl{Among causal laws of the restricted form $u_{1,k}=-\widehat K_k x_k$ (no feedforward),
% the CF law $u_{1,k}^{\mathrm{CF}}=-K_k x_k - L_k w_k$ uniquely minimizes $J_1$, and
% for any $\{\widehat K_k\}$,}
% \[
% J_1(\{-\widehat K_k x_k\}) - J_1(\mathrm{CF})
% = \sum_{k=0}^{N-1} \|\,H_k^{-\tfrac12}B_1^\top P_{k+1} w_k\,\|_2^2 \;\ge 0,
% \]
% with strict inequality whenever $w_k\not\equiv 0$ in a controllable–costate direction.

\subsection{Specialization: Slowly-Varying Error Dynamics}
When the FNE command $u_{2,k}^\star$ varies slowly, the drive term becomes negligible and $w_{k+1}=\alpha w_k$. The CF gains then admit closed forms in terms of the block Riccati sequence.

\begin{corollary}
\label{cor:slow-varying-alpha}
Assume the lag model \eqref{eq:w-lag} with $0<\alpha<1$ and that Player~2’s Nash command varies slowly so that
\begin{equation}\label{eq:slow-drive}
u^{\star}_{2,k+1}\approx u^{\star}_{2,k}
\end{equation}
hence the drive is negligible and $w_{k+1} = \alpha\,w_k$. Then the optimal gains $K_k$ and $L_k$ defined in Theorem~\ref{thm:lag-optimal} will be obtained as
\begin{equation}
\begin{aligned}
&K_k \;=\; H_k^{-1} B_1^\top P_{k+1}^{xx} A_{\mathrm{eff},k}, \\
&L_k \;=\; H_k^{-1} B_1^\top \big(P_{k+1}^{xx} + \alpha\,P_{k+1}^{xw}\big), \label{eq:special case k}
\end{aligned}
\end{equation}
where $H_k \;:=\; R_1 + B_1^\top P_{k+1}^{xx} B_1 \ \succ 0$; $P_{k}^{xx}$ satisfies the standard finite-horizon Riccati recursion
\begin{equation}\label{eq:Riccati-Pxx}
\begin{aligned}
&P_N^{xx}=Q_{1,N},\\
&\begin{aligned}
P_k^{xx}=Q_1 + &A_{\mathrm{eff},k}^\top P_{k+1}^{xx} A_{\mathrm{eff},k}\\
&- A_{\mathrm{eff},k}^\top P_{k+1}^{xx} B_1 H_k^{-1} B_1^\top P_{k+1}^{xx} A_{\mathrm{eff},k},
\end{aligned}
\end{aligned}
\end{equation}
and the cross block $P_k^{xw}$ evolves linearly as
\begin{equation}\label{eq:Riccati-Pxw}
\begin{aligned}
&P_N^{xw}=0,\\
&\begin{aligned}
P_k^{xw} =\; &A_{\mathrm{eff},k}^\top P_{k+1}^{xx}
+ \alpha\,A_{\mathrm{eff},k}^\top P_{k+1}^{xw}\\
&- A_{\mathrm{eff},k}^\top P_{k+1}^{xx} B_1 H_k^{-1} B_1^\top \big(P_{k+1}^{xx}+\alpha\,P_{k+1}^{xw}\big).
\end{aligned}
\end{aligned}
\end{equation}
\end{corollary}

\begin{proof}
The assumption \eqref{eq:slow-drive} implies $w_{k+1}=\alpha w_k$. In augmented coordinates
$z_k=[x_k^\top \; w_k^\top]^\top$ this gives
\[
\bar A_k=\begin{bmatrix}A_{\mathrm{eff},k}& I\\ 0 & \alpha I\end{bmatrix},\qquad
\bar B_k=\begin{bmatrix}B_1\\ 0\end{bmatrix}.
\]
By Theorem~\ref{thm:lag-optimal}, the optimal control \eqref{eq:ukstar} can be simplified by partition $\bar P_{k+1}=\begin{bmatrix}P_{k+1}^{xx}&P_{k+1}^{xw}\\ P_{k+1}^{wx}&P_{k+1}^{ww}\end{bmatrix}$ and computing the terms
\begin{equation*}
\bar B_k^\top \bar P_{k+1}\bar B_k \;=\; B_1^\top P_{k+1}^{xx} B_1,
\end{equation*}
\begin{equation*}
\bar B_k^\top \bar P_{k+1}\bar A_k
= B_1^\top\!\big[\,P_{k+1}^{xx}A_{\mathrm{eff},k}\ \ \ P_{k+1}^{xx}+\alpha P_{k+1}^{xw}\,\big].
\end{equation*}
Defining $H_k := R_1 + B_1^\top P_{k+1}^{xx} B_1 \succ 0$, 
\begin{equation}
\begin{aligned}
u_{1,k}^\star
&= -\,H_k^{-1}\,B_1^\top\!\big[\,P_{k+1}^{xx}A_{\mathrm{eff},k}\ \ \ P_{k+1}^{xx}+\alpha P_{k+1}^{xw}\,\big]
\begin{bmatrix}x_k\\ w_k\end{bmatrix}\\[2pt]
&\;\begin{aligned}
= &-\,\underbrace{H_k^{-1}B_1^\top P_{k+1}^{xx}A_{\mathrm{eff},k}}_{\displaystyle K_k}\,x_k\\
&\;-\;\underbrace{H_k^{-1}B_1^\top\big(P_{k+1}^{xx}+\alpha P_{k+1}^{xw}\big)}_{\displaystyle L_k}\,w_k,
\end{aligned}
\end{aligned}
\end{equation}
which yields \eqref{eq:special case k}. Taking the $(1,1)$ and $(1,2)$ blocks of the Riccati recursion
\eqref{eq:riccati-aug} with the above $\bar A_k,\bar B_k$ gives the stated recursions
\eqref{eq:Riccati-Pxx}–\eqref{eq:Riccati-Pxw} and $P_N^{xw}=0$ follows from $\bar Q_N$ being block-diagonal.
\end{proof}

Comparing the Riccati recursion $P^{xx}_k$ in \eqref{eq:Riccati-Pxx}, and the gain formulation $K_k$ in \eqref{eq:special case k} with the ones in Definition~\ref{def: definition of Nash} reveals that the state gain $K_k$ obtained under the assumption of Corollary~\ref{cor:slow-varying-alpha} is equal to that of the FNE.  


It is worth noting that \eqref{eq:Riccati-Pxx} coincides with the standard finite-horizon Riccati recursion for the disturbance-free system. In particular, $K_k$ matches the usual state-feedback gain derived from $P_k^{xx}$, while $L_k$ adds a causal feedforward term shaped by both $P_{k+1}^{xx}$ and the cross block $P_{k+1}^{xw}$ scaled by $\alpha$. The resulting law $u_{1,k}^\star=-K_k x_k - L_k w_k$ thus has the familiar structure of disturbance-accommodation or preview control: the controller reacts not only to the state but also to the measured disturbance. Under the lag generator \eqref{eq:w-lag}, the augmented Riccati \eqref{eq:riccati-aug} introduces off-diagonal couplings between $x$ and $w$; these generate a preview component in $L_k$ that anticipates how current deviations propagate through the lag into future stages. Hence \eqref{eq:gains-closed} recovers the classical disturbance-accommodation principle and extends it to physically grounded, time-varying deviation dynamics.


% Let the control action of Player~1 be in modified feedback Nash strategy form as defined in Equation (\ref{eq: optimal strategy for P1 in disturbance situation}). This makes an analogy with the standard LQR problem with optimal offline policy addressed in \cite{goel2022power}. The game-theoretic nature of the problem is embedded in $A_{\mathrm{eff}}$ by including Player~2's dynamics.
    % \hl{However, for the sake of simplicity, the knowledge of future disturbances is not considered in this paper}. 
    % For the linear dynamic system (\ref{eq: modified dynamics}) and the quadratic cost function (\ref{eq: LQ cost function}) for Player~1 ($i=1$), solving the dynamic programming approach \cite{bellman1966dynamic} leads to solving the following \hl{DRE}:

% \begin{corollary}
% \label{cor:slow-varying-alpha}
% Assume the lag model \eqref{eq:w-lag} with $0<\alpha<1$ and that Player~2’s Nash command varies slowly so that
% \begin{equation}\label{eq:slow-drive}
% u^{\star}_{2,k+1}\approx u^{\star}_{2,k}\quad\Longrightarrow\quad
% K^{\star}_{2,k+1}\approx K^{\star}_{2,k},
% \end{equation}
% hence the drive is negligible and $w_{k+1} = \alpha\,w_k$. Then the augmented dynamics in Theorem~\ref{thm:lag-optimal} are accurately approximated by
% \begin{equation}\label{eq:ABQ-slow}
% \bar A_k \;\approx\; \begin{bmatrix} A_{\mathrm{eff},k} & I\\[2pt] 0 & \alpha I \end{bmatrix},
% \qquad
% \bar B_k \;\approx\; \begin{bmatrix} B_1\\[2pt] 0 \end{bmatrix},
% \end{equation}
% Let $\bar P_k=\begin{bmatrix} P_k^{xx} & P_k^{xw}\\ P_k^{wx} & P_k^{ww}\end{bmatrix}$ be the Riccati sequence for \eqref{eq:ABQ-slow}. Then:
% \begin{equation}
% \begin{aligned}
% &H_k \;:=\; R_1 + B_1^\top P_{k+1}^{xx} B_1 \ \succ 0, \label{eq:Hk-slow}
% \end{aligned}
% \end{equation}
% \begin{equation}
% \begin{aligned}
% &K_k \;=\; H_k^{-1} B_1^\top P_{k+1}^{xx} A_{\mathrm{eff},k}, \\
% &L_k \;=\; H_k^{-1} B_1^\top \big(P_{k+1}^{xx} + \alpha\,P_{k+1}^{xw}\big), \label{eq:gains-slow}
% \end{aligned}
% \end{equation}
% where $P_{k}^{xx}$ satisfies the standard finite-horizon Riccati recursion
% \begin{equation}\label{eq:Riccati-Pxx}
% \begin{aligned}
% &P_N^{xx}=Q_{1,N},\\
% &\begin{aligned}
% P_k^{xx}=Q_1 + &A_{\mathrm{eff},k}^\top P_{k+1}^{xx} A_{\mathrm{eff},k}\\
% &- A_{\mathrm{eff},k}^\top P_{k+1}^{xx} B_1 H_k^{-1} B_1^\top P_{k+1}^{xx} A_{\mathrm{eff},k},
% \end{aligned}
% \end{aligned}
% \end{equation}
% and the cross block $P_k^{xw}$ evolves linearly as
% \begin{equation}\label{eq:Riccati-Pxw}
% \begin{aligned}
% P_k^{xw} =\; &A_{\mathrm{eff},k}^\top P_{k+1}^{xx}
% + \alpha\,A_{\mathrm{eff},k}^\top P_{k+1}^{xw}\\
% &- A_{\mathrm{eff},k}^\top P_{k+1}^{xx} B_1 H_k^{-1} B_1^\top \big(P_{k+1}^{xx}+\alpha\,P_{k+1}^{xw}\big),
% \\
% P_N^{xw}=0.
% \end{aligned}
% \end{equation}
% In particular, $K_k$ is exactly the state-feedback from the standard Riccati in $P_k^{xx}$, while $L_k$ adds a preview/feedforward term that depends on both $P_{k+1}^{xx}$ and the cross block $P_{k+1}^{xw}$ scaled by~$\alpha$.
% \end{corollary}

\begin{remark}
To sanity-check the design, we examine the two extreme points of the lag parameter 
$\alpha$ in the implemented-command model \eqref{eq:u2-lag}. 
When $\alpha\to 0$, the actuator tracks its Nash command essentially instantaneously, so 
$u_{2,k+1}\approx u^{\star}_{2,k}$ and the disturbance recursion reduces to 
$w_{k+1}=-B_2(u^{\star}_{2,k+1}-u^{\star}_{2,k})$. In steady segments the disturbance vanishes, 
and the compensated feedback coincides with the classical disturbance-accommodation (DBSC) law 
that uses $w_k$ as an immediate feedforward signal. At the other extreme, when $\alpha\to 1$ and 
the drive term is negligible, the disturbance behaves like a nearly constant bias with 
$w_{k+1}\approx w_k$, so the compensation acts as preview cancellation of a persistent input. 
Both cases are contained as limits of the proposed lag model and show that the compensated law 
weakly dominates any $K$-only feedback whenever a measurable deviation is present.
\end{remark}

% \begin{remark}
% \label{rem:alpha-extremes}
% It is instructive to examine two extreme points of the lag parameter $\alpha$. When $\alpha\to 0$, the actuator tracks its Nash command almost instantaneously, so execution errors do not persist. In this limit, the disturbance vanishes during steady segments, and the compensated feedback collapses to the instantaneous preview/DBSC case. At the other extreme, when $\alpha\to 1$ and the drive term is negligible, the disturbance behaves like a constant bias with 
% $w_{k+1} \approx w_k$. The controller then reduces to a preview cancellation of a nearly constant input. These two limits highlight that the compensated law consistently weakly dominates any $K_k$-only policy, with equality only if the disturbance is identically zero or acts along uncontrollable/cost-orthogonal directions.

% (i) If $\alpha=1$ and the drive vanishes exactly ($u^{\star}_{2,k+1}=u^{\star}_{2,k}$), then $w_{k+1}=w_k$ and the augmented pair reduces further to
% $\bar A_k=\begin{bmatrix}A_{\mathrm{eff},k}&I\\ 0&I\end{bmatrix}$, $\bar B_k=\begin{bmatrix}B_1\\ 0\end{bmatrix}$; the optimal gains collapse to the classical DBSC/preview form
% $K_k=H_k^{-1}B_1^\top P_{k+1}^{xx}A_{\mathrm{eff},k}$, $L_k=H_k^{-1}B_1^\top P_{k+1}^{xx}$ with $P^{xx}_k$ from \eqref{eq:Riccati-Pxx}. \\
% (ii) If $\alpha=0$ (very fast actuator), then $u_{2,k+1}=u^{\star}_{2,k}$ and the disturbance does not persist: $w_{k+1}=-B_2(u^{\star}_{2,k+1}-u^{\star}_{2,k})$. In steady stretches ($u^{\star}_{2,k+1}=u^{\star}_{2,k}$) this gives $w_{k+1}=0$, recovering the instantaneous-disturbance case.
% (i) As $\alpha\to 0$ (very fast implementation), the lag vanishes and the CF law converges to the instantaneous-preview (DBSC) case. 
% (ii) As $\alpha\to 1$ with negligible drive, $w_{k+1}\approx w_k$ (persistent bias); CF reduces to preview cancellation of a constant disturbance. In both cases, CF weakly dominates any $K$-only controller; equality occurs only if $w_k\equiv 0$ or along uncontrollable/cost-orthogonal directions.

% (i) If $\alpha\to 0$ (very fast implementation), \eqref{eq:augAB} reduces to the instantaneous-measurement case, and $H_k$ collapses to the classical DBSC feedforward gain $L_k$. 
% (ii) If $u^{\star}_{2,k}$ is (piecewise) constant over the horizon so that $u^{\star}_{2,k+1}\approx u^{\star}_{2,k}$, then $w_{k+1}\approx \alpha w_k$ and \eqref{eq:gains-aug} specializes to the simpler AR(1) preview controller with feedforward proportional to the cross-block of $\bar P_{k+1}$.
% \end{remark}


% \begin{proof}
% Equation~(\ref{eq: augmented LQR}) represents a standard finite-horizon discrete-time LQR problem for the augmented system. Its optimal control input is obtained from the augmented state-feedback policy
% \begin{equation*}
%     u_{1,k} = \bar{K}_k \bar{x}_k,
% \end{equation*}
% where the optimal gain $\bar{K}_k$ is computed offline via the recursion \cite{stengel1994optimal},
% \begin{equation}
% \label{eq: optimal augmented gain}
%     \bar{K}_k = (R_1 + \bar{B}^\top \bar{P}_{k+1} \bar{B})^{-1} \bar{B}^\top \bar{P}_{k+1} \bar{A},
% \end{equation}
% with $\bar{P}_{k+1}$ satisfying the following DRE:
% \begin{equation}
%     \label{eq: Riccati for standard LQR}
%     \begin{aligned}
%         \bar{P}_k &= \bar{Q} + \bar{A}^\top \bar{P}_{k+1} \bar{A}  \\&- \bar{A}^\top \bar{P}_{k+1} \bar{B} (R_1 + \bar{B}^\top \bar{P}_{k+1} \bar{B})^{-1} \bar{B}^\top \bar{P}_{k+1} \bar{A}.
%     \end{aligned}
% \end{equation}

% Assume the solution \( \bar{P}_k \) admits the block decomposition:
% \begin{equation*}
%     \bar{P}_k = \begin{bmatrix} P_k^{xx} & P_k^{xw} \\ P_k^{wx} & P_k^{ww} \end{bmatrix}.
% \end{equation*}

% Expanding the optimal gain $\bar{K}_k$ defined in (\ref{eq: optimal augmented gain}) using this block structure for $\bar{P}_{k+1}$, we obtain:
% \begin{align*}
%     \bar{K}_k &= \left(R_1+\begin{bmatrix} B_1 \\ 0 \end{bmatrix}^\top\begin{bmatrix} P_k^{xx} & P_k^{xw} \\ P_k^{wx} & P_k^{ww} \end{bmatrix}\begin{bmatrix} B_1 \\ 0 \end{bmatrix}\right)^{-1}\\
%     & \times \begin{bmatrix} B_1 \\ 0 \end{bmatrix}^\top\begin{bmatrix} P_k^{xx} & P_k^{xw} \\ P_k^{wx} & P_k^{ww} \end{bmatrix} \begin{bmatrix} A_{\mathrm{eff}} & I \\ 0 & 0 \end{bmatrix}.
% \end{align*}

% This expression results in a $2\times 1$ block matrix, which can be partitioned as
% \begin{equation*}
%     \bar{K}_k=\begin{bmatrix} K^\star_{1,k} & L^\star_k \end{bmatrix},
% \end{equation*}
% % \begin{align*}
% %     &\bar{K}_k = \nonumber \\& \begin{bmatrix} (R_1 + B_1^\top P_{k+1}^{xx}B_1)^{-1}B_1^\top P_{k+1}^{xx}A_{\mathrm{eff}} & (R_1 + B_1^\top P_{k+1}^{xx}B_1)^{-1}B_1^\top P_{k+1}^{xx} \end{bmatrix}
% % \end{align*}
% with the optimal gains,
% \begin{align*}
%     K^\star_{1,k} &= (R_1 + B_1^\top P_{k+1}^{xx} B_1)^{-1} B_1^\top P_{k+1}^{xx} A_{\mathrm{eff}}, \\
%     L^\star_k &= (R_1 + B_1^\top P_{k+1}^{xx} B_1)^{-1} B_1^\top P_{k+1}^{xx}.
% \end{align*}

% To compute $P_{k+1}^{xx}$, we substitute the block decomposition $\bar{P}_{k+1}$ into the DRE (\ref{eq: Riccati for standard LQR}). 

% After some algebraic manipulation, the block component $P_{k+1}^{xx}$ satisfies:
% \begin{align*}
%     P_k^{xx} &= Q_1 + A_{\mathrm{eff}}^\top P^{xx}_{k+1} A_{\mathrm{eff}} \nonumber \\&- A_{\mathrm{eff}}^\top P^{xx}_{k+1} B_1 (R_1 + B_1^\top P^{xx}_{k+1} B_1)^{-1} B_1^\top P^{xx}_{k+1} A_{\mathrm{eff}}.
% \end{align*}
% This expression is the standard finite-horizon DRE for the LQR problem applied to the original, non-augmented, disturbance-free system dynamics, with the solution denoted by $P_k$ as in \eqref{eq: DRE for P1}. This concludes the proof.
% \end{proof}

% \begin{proof}
% returns the optimal policy of,
%     \begin{equation*}
%         u_{1,k} = -\bar{K}_k \bar{x}_k,
%     \end{equation*}
% where
%     \begin{equation}
%         \label{eq: optimal gain in standard LQR}
%         \bar{K}_k = (R_1 + \bar{B}^\top \bar{P}_{k+1} \bar{B})^{-1} \bar{B}^\top \bar{P}_{k+1} \bar{A}.
%     \end{equation}
% and $\bar{P}_{k+1}$ is the solution of the Riccati recursion,
%     \begin{align}
%         \label{eq: Riccati for standard LQR}
%         \bar{P}_k &= \bar{Q} + \bar{A}^\top \bar{P}_{k+1} \bar{A} \nonumber \\&- \bar{A}^\top \bar{P}_{k+1} \bar{B} (R_1 + \bar{B}^\top \bar{P}_{k+1} \bar{B})^{-1} \bar{B}^\top \bar{P}_{k+1} \bar{A},
%     \end{align}
% with the terminal condition $\bar{P}_N = \bar{Q}_N$.

% Expanding the optimal gain obtained in Equation~(\ref{eq: optimal gain in standard LQR}) results in the following expression,
% \begin{align*}
%     \bar{K}_k = \left(R_1+\begin{bmatrix} B_1 \\ 0 \end{bmatrix}^\top\bar{P}_{k+1}\begin{bmatrix} B_1 \\ 0 \end{bmatrix}\right)^{-1}\begin{bmatrix} B_1 \\ 0 \end{bmatrix}^\top\bar{P}_{k+1} \begin{bmatrix} A_{\mathrm{eff}} & I \\ 0 & 0 \end{bmatrix}.
% \end{align*}

% Let assume $P_{k+1}$ is a $2\times 2$ block matrix as follows,
% \[
% \bar{P}_{k+1} = \begin{bmatrix} P_{k+1}^{xx} & P_{k+1}^{xw} \\ P_{k+1}^{wx} & P_{k+1}^{ww} \end{bmatrix}.
% \]

% Substituting the block matrix in the expanded expression results in the following matrix after some manipulations,
% \begin{align*}
%     &\bar{K}_k = \nonumber \\& \begin{bmatrix} (R_1 + B_1^\top P_{k+1}^{xx}B_1)^{-1}B_1^\top P_{k+1}^{xx}A_{\mathrm{eff}} & (R_1 + B_1^\top P_{k+1}^{xx}B_1)^{-1}B_1^\top P_{k+1}^{xx} \end{bmatrix}
% \end{align*}

% Splitting $\bar{K}_k$ into $\bar{K}_k = \begin{bmatrix} K_{1,k}^\star & L_k^\star \end{bmatrix}$ gives the optimal matrices. 

% To obtain the block matrix $P_{k+1}^{xx}$, it is required to expand the Riccati recursion given in Equation~(\ref{eq: Riccati for standard LQR}).
% % \begin{align}
% %     \bar{K}_k &= \left(R_1+\begin{bmatrix} B_1^\top & 0 \end{bmatrix}
% %     \begin{bmatrix} P_{k+1}^{xx} & P_{k+1}^{xw} \\ P_{k+1}^{wx} & P_{k+1}^{ww} \end{bmatrix}
% %     \begin{bmatrix} B_1 \\ 0 \end{bmatrix}\right)^{-1} \notag \\
% %     &\quad \times \begin{bmatrix} B_1^\top & 0 \end{bmatrix}
% %     \begin{bmatrix} P_{k+1}^{xx} & P_{k+1}^{xw} \\ P_{k+1}^{wx} & P_{k+1}^{ww} \end{bmatrix} 
% %     \begin{bmatrix} A_{\mathrm{eff}} & I_{\bar{n}\times \bar{n}} \\ 0 & 0 \end{bmatrix}
% % \end{align}
% Rewriting the Riccati equation,
% \begin{align*}
%         \label{eq: Riccati for standard LQR}
%         \bar{P}_k = \bar{Q} + \bar{A}^\top \bar{P}_{k+1} \bar{A} \nonumber - \bar{A}^\top \bar{P}_{k+1} \bar{B} \bar{K}_k,
%     \end{align*}
% results in,
% \begin{align*}
%     \begin{bmatrix} P_{k}^{xx} & P_{k}^{xw} \\ P_{k}^{wx} & P_{k}^{ww} \end{bmatrix} 
%     &= \begin{bmatrix} Q_1 & 0 \\ 0 & 0 \end{bmatrix} \notag \\
%     & + \begin{bmatrix} A_{\mathrm{eff}}^\top & 0 \\ I & 0 \end{bmatrix}
%     \begin{bmatrix} P_{k+1}^{xx} & P_{k+1}^{xw} \\ P_{k+1}^{wx} & P_{k+1}^{ww} \end{bmatrix}
%     \begin{bmatrix} A_{\mathrm{eff}} & I \\ 0 & 0 \end{bmatrix} \notag \\
%     & + \begin{bmatrix} A_{\mathrm{eff}}^\top & 0 \\ I & 0 \end{bmatrix}
%     \begin{bmatrix} P_{k+1}^{xx} & P_{k+1}^{xw} \\ P_{k+1}^{wx} & P_{k+1}^{ww} \end{bmatrix}
%     \begin{bmatrix} B_1 \\ 0 \end{bmatrix} \notag \\
%     & \times \left(R_1 + \begin{bmatrix} B_1^\top & 0 \end{bmatrix}
%     \begin{bmatrix} P_{k+1}^{xx} & P_{k+1}^{xw} \\ P_{k+1}^{wx} & P_{k+1}^{ww} \end{bmatrix}
%     \begin{bmatrix} B_1 \\ 0 \end{bmatrix} \right)^{-1} \notag \\
%     & \times \begin{bmatrix} B_1^\top & 0 \end{bmatrix}
%     \begin{bmatrix} P_{k+1}^{xx} & P_{k+1}^{xw} \\ P_{k+1}^{wx} & P_{k+1}^{ww} \end{bmatrix} 
%     \begin{bmatrix} A_{\mathrm{eff}} & I \\ 0 & 0 \end{bmatrix}
% \end{align*}

% % \begin{align*}
% %         \begin{bmatrix} P_{k}^{xx} & P_{k}^{xw} \\ P_{k}^{wx} & P_{k}^{ww} \end{bmatrix} & = \begin{bmatrix} Q_1 & 0 \\ 0 & 0 \end{bmatrix} + \begin{bmatrix} A_{\mathrm{eff}}^\top & 0 \\ I & 0 \end{bmatrix}\begin{bmatrix} P_{k+1}^{xx} & P_{k+1}^{xw} \\ P_{k+1}^{wx} & P_{k+1}^{ww} \end{bmatrix}\begin{bmatrix} A_{\mathrm{eff}} & I \\ 0 & 0 \end{bmatrix}\\
% %         & + \begin{bmatrix} A_{\mathrm{eff}}^\top & 0 \\ I & 0 \end{bmatrix}\begin{bmatrix} P_{k+1}^{xx} & P_{k+1}^{xw} \\ P_{k+1}^{wx} & P_{k+1}^{ww} \end{bmatrix}\begin{bmatrix} B_1 \\ 0 \end{bmatrix}\left(R_1+\begin{bmatrix} B_1^\top & 0 \end{bmatrix}\begin{bmatrix} P_{k+1}^{xx} & P_{k+1}^{xw} \\ P_{k+1}^{wx} & P_{k+1}^{ww} \end{bmatrix}\begin{bmatrix} B_1 \\ 0 \end{bmatrix}\right)^{-1}\begin{bmatrix} B_1^\top & 0 \end{bmatrix}\begin{bmatrix} P_{k+1}^{xx} & P_{k+1}^{xw} \\ P_{k+1}^{wx} & P_{k+1}^{ww} \end{bmatrix} \begin{bmatrix} A_{\mathrm{eff}} & I \\ 0 & 0 \end{bmatrix}
% % \end{align*}

% After some manipulation, $P_{k+1}^{xx}$ is obtained from,
% \begin{align*}
%     % P_k^{xx} = Q_1 +  A_{\mathrm{eff}} \\
%     P_k^{xx} &= Q_1 + A_{\mathrm{eff}}^\top P^{xx}_{k+1} A_{\mathrm{eff}} \nonumber \\&- A_{\mathrm{eff}}^\top P^{xx}_{k+1} B_1 (R_1 + B_1^\top P^{xx}_{k+1} B_1)^{-1} B_1^\top P^{xx}_{k+1} A_{\mathrm{eff}},
% \end{align*}
% which is the DRE for non-augmmented non-disturbed system and can be denoted by $P_{k+1}$.

%     % Substituting \( u_k \) into the original system:
%     % \begin{equation}
%     %     x_{k+1} = A x_k + B u_k + w_k.
%     % \end{equation}
%     % \begin{equation}
%     %     x_{k+1} = A x_k + B(-K_k x_k + L_k w_k) + w_k.
%     % \end{equation}
%     % \begin{equation}
%     %     x_{k+1} = (A - B K_k) x_k + (I + B L_k) w_k.
%     % \end{equation}

%     % The disturbance effect is governed by:
%     % \begin{equation}
%     %     \| I + B L_k \|.
%     % \end{equation}

%     % If \( \| I + B L_k \| \ll 1 \), the effect of \( w_k \) is significantly reduced.


    
%     Let the control action of Player~1 be in modified feedback Nash strategy form as defined in Equation (\ref{eq: optimal strategy for P1 in disturbance situation}). This makes an analogy with the standard LQR problem with optimal offline policy addressed in \cite{goel2022power}. The game-theoretic nature of the problem is embedded in $A_{\mathrm{eff}}$ by including Player~2's dynamics.
%     % \hl{However, for the sake of simplicity, the knowledge of future disturbances is not considered in this paper}. 
%     For the linear dynamic system (\ref{eq: modified dynamics}) and the quadratic cost function (\ref{eq: LQ cost function}) for Player~1 ($i=1$), solving the dynamic programming approach \cite{bellman1966dynamic} leads to solving the following \hl{DRE}:
%     \begin{align}
%         P_k &= Q_1 + A_{\mathrm{eff}}^\top P_{k+1} A_{\mathrm{eff}} \nonumber \\
%       & \;-\; A_{\mathrm{eff}}^\top P_{k+1} B_1 
%       \bigl(R_1 + B_1^\top P_{k+1} B_1\bigr)^{-1} 
%       B_1^\top P_{k+1} A_{\mathrm{eff}}, \nonumber \\
%       P_N &= Q_{1,N}.
%     \end{align}
%     The solution of \hl{DRE} gives the optimal control gains (\ref{eq: optimal gains}) for the control policy (\ref{eq: optimal strategy for P1 in disturbance situation}) and ends the proof.
% \end{proof}

% ******************************************************


% \begin{remark}[Estimation and certainty-equivalence]
% If Player~1 has only a causal estimate $\hat w_k$ of $w_k$, the same gains $(K_k,L_k)$ apply (certainty equivalence) when $\hat w_k$ is produced by a linear observer consistent with \eqref{eq:augAB}. The controller uses $u_{1,k}=-K_k x_k - H_k \hat w_k$.
% \end{remark}


% \begin{definition}
% \label{def:value-function}
% For $k=N$, define
% \begin{equation}
% V_N(x) := x^\top Q_{1,N} x .
% \end{equation}
% For $k=N-1,\dots,0$, define the cost-to-go at time $k$ from state $x$ as
% \begin{equation}
% \label{eq:Vk-def}
% \begin{split}
% V_k(x):=\min_{\{u_{1,\ell}\}_{\ell=k}^{N-1}}
% \Big\{ \sum_{\ell=k}^{N-1}\!\big(x_\ell^\top Q_1 x_\ell + &u_{1,\ell}^\top R_1 u_{1,\ell}\big)\\
% &\;+\; x_N^\top Q_{1,N} x_N \Big\}.
% \end{split}
% \end{equation}

% subject to \eqref{eq:plant-general}, with the information pattern of Assumption~\ref{ass:info-and-weights}. In the linear–quadratic setting, the Bellman recursion implies
% \begin{equation}
% \label{eq:Vk-quadratic}
% V_k(x) = x^\top P_k x + c_k,
% \end{equation}
% for some $P_k\succeq 0$ and scalar $c_k$, with $P_N=Q_{1,N}$.
% \end{definition}


% \begin{remark}
% In \eqref{eq:Vk-def} and \eqref{eq:Vk-quadratic}, the argument $x$ denotes the \emph{state at time $k$}. Inside the Bellman recursion \eqref{eq:bellman}, the only future quantity that enters is $V_{k+1}(x_{k+1})$, which \emph{already} summarizes the optimal cost over all $\{x_\ell,u_{1,\ell}\}_{\ell\ge k+1}$ starting from the next state $x_{k+1}$.
% \end{remark}


% \begin{proof}
% We use dynamic programming with the Bellman recursion. Define the value function $V_k(x)$ as
% \begin{equation}
% \label{eq:Vk-def}
%  \min_{\{u_{1,\ell}\}_{\ell=k}^{N-1}} 
% \left\{ \sum_{\ell=k}^{N-1}\big(x_\ell^\top Q_1 x_\ell + u_{1,\ell}^\top R_1 u_{1,\ell}\big) + x_N^\top Q_{1,N}x_N \right\},
% \end{equation}
% subject to \eqref{eq:plant-general} and the information pattern of Assumption~\ref{ass:info-and-weights}. The terminal condition is $V_N(x)=x^\top Q_{1,N}x$.

% Assume $V_{k+1}(x)=x^\top P_{k+1}x + c_{k+1}$ for some $P_{k+1}\succeq0$, and  $c_{k+1}$.

% At time $k$, given $(x_k,w_k)$,
% \begin{equation}
% \label{eq:bellman}
% V_k(x_k)=\min_{u_{1,k}}\Big\{ x_k^\top Q_1 x_k + u_{1,k}^\top R_1 u_{1,k} + V_{k+1}(x_{k+1}) \Big\}.
% \end{equation}
% Substituting $x_{k+1}$ from \eqref{eq:plant-general} into $V_{k+1}(x_{k+1})$, expanding it, and collecting terms in $u_{1,k}$ yields the strictly convex quadratic
% \begin{equation}
% \label{eq:Qk}
% \begin{aligned}
% \mathcal Q_k(x_k,u_{1,k})
% &= u_{1,k}^\top H_k u_{1,k}
% + 2\,u_{1,k}^\top B_1^\top P_{k+1}\big(A_{\text{eff},k}x_k+w_k\big)\\
% &+ (A_{\text{eff},k}x_k+w_k)^\top P_{k+1}(A_{\text{eff},k}x_k+w_k)\\
% &+ x_k^\top Q_1 x_k + c_{k+1}.
% \end{aligned}
% \end{equation}
% Hence, the unique minimizer is
% \begin{equation}
% \label{eq:KL-sep}
% u_{1,k}^\star \;=\; -\underbrace{H_k^{-1}B_1^\top P_{k+1}A_{\text{eff},k}}_{K_k}\,x_k
% \;-\;\underbrace{H_k^{-1}B_1^\top P_{k+1}}_{L_k}\,w_k.
% \end{equation}
% Thus \eqref{eq:KL-opt} holds.
% Substituting $u_{1,k}^\star$ into \eqref{eq:bellman} and matching the quadratic term in $x_k$ gives the Riccati update \eqref{eq:riccati-Aeff}.

% For any alternative $\tilde u_{1,k}$,
% \begin{equation}
% \label{eq:one-step-gap}
% \begin{aligned}
% \mathcal Q_k(x_k,\tilde u_{1,k}) &- \mathcal Q_k(x_k,u_{1,k}^\star)\\
% &= (\tilde u_{1,k}-u_{1,k}^\star)^\top H_k (\tilde u_{1,k}-u_{1,k}^\star) \;\ge\; 0.
% \end{aligned}
% \end{equation}
% Summing \eqref{eq:one-step-gap} over $k=0,\dots,N-1$ telescopes the value functions and yields \eqref{eq:gap-sos}. Choosing $(\tilde K_k,\tilde L_k)=(K^\star_{1,k},0)$ gives \eqref{eq:dominance-stick} and completes this proof.
% \end{proof}


% % ****************************************************
% \section{Best Response with Measured Deviations (General Case)}
% \label{sec:best-response-general}

% Let $w_k := B_2\,\Delta u_{2,k}$ denote the (measured/estimable) execution error of Player~2 at time $k$.
% Fix Player~2 at its FNE feedback $u^\star_{2,k}=-K^\star_{2,k}x_k$ and define the effective matrices
% \begin{equation}
% \label{eq:Aeff-def}
% A_{\text{eff},k}:=A-B_2K^\star_{2,k}, \qquad k=0,\dots,N-1.
% \end{equation}
% Then the state evolves according to
% \begin{equation}
% \label{eq:general-plant}
% x_{k+1}=A_{\text{eff},k}\,x_k + B_1\,u_{1,k} + w_k .
% \end{equation}

% \begin{assumption}[Information available to Player~1]
% \label{ass:measured-wk}
% At time $k$, Player~1 knows the current state $x_k$ and the current deviation input $w_k$ (but not future $w_{k+i}$, $i\ge1$).
% Matrices $Q_1\succeq 0$, $Q_{1,N}\succeq 0$, and $R_1\succ 0$.
% \end{assumption}

% \begin{theorem}[Best response with instantaneous measurement; dominance over stick-with-Nash]
% \label{thm:best-response-general}
% Under Assumptions~\ref{ass:memoryless-feedback} and \ref{ass:measured-wk}, consider the causal affine policy class
% \begin{equation}
% \label{eq:policy-class}
% u_{1,k} = -K_k x_k - L_k w_k, \qquad k=0,\dots,N-1,
% \end{equation}
% with time-varying gains $\{K_k,L_k\}$.
% Define $P_N:=Q_{1,N}$ and the backward Riccati recursion
% \begin{equation}
% \label{eq:riccati-Aeff}
% \begin{aligned}
%     P_k \;&=\; Q_1 + A_{\text{eff},k}^\top P_{k+1} A_{\text{eff},k}\\
%     &\; - A_{\text{eff},k}^\top P_{k+1} B_1 \big(R_1 + B_1^\top P_{k+1} B_1\big)^{-1} B_1^\top P_{k+1} A_{\text{eff},k},
% \end{aligned}
% \end{equation}
% for $k=N-1,\dots,0$. Let $H_k:=R_1+B_1^\top P_{k+1}B_1\succ 0$.
% Then the gains
% \begin{equation}
% \label{eq:KL-opt}
% K_k \;=\; H_k^{-1} B_1^\top P_{k+1} A_{\text{eff},k}, 
% \qquad
% L_k \;=\; H_k^{-1} B_1^\top P_{k+1},
% \end{equation}
% minimize Player~1's cost $J_1$ over all causal laws of the form \eqref{eq:policy-class} for any disturbance sequence $\{w_k\}$, i.e.,
% \begin{equation}
% \label{eq:dominance-any}
% J_1(\{K_k,L_k\}) \;\le\; J_1(\{\tilde K_k,\tilde L_k\}) \quad \text{for all causal } \{\tilde K_k,\tilde L_k\}.
% \end{equation}
% In particular, the stick-with-Nash baseline $u_{1,k}=-K^\star_{1,k}x_k$ corresponds to $(\tilde K_k,\tilde L_k)=(K^\star_{1,k},0)$, hence
% \begin{equation}
% \label{eq:dominance-stick}
% J_1(\{K_k,L_k\}) \;\le\; J_1(\{K^\star_{1,k},0\}) \quad \text{for every sequence } \{w_k\}.
% \end{equation}
% Moreover, for any causal $\{\tilde K_k,\tilde L_k\}$ the cost gap admits the nonnegative quadratic decomposition
% \begin{equation}
% \label{eq:gap-sum-of-squares}
% \begin{aligned}
%     J_1(\{\tilde K_k,\tilde L_k\}) &- J_1(\{K_k,L_k\})\\
%     \;&=\; \sum_{k=0}^{N-1} \big(\tilde u_{1,k}-u_{1,k}^\star\big)^\top H_k \big(\tilde u_{1,k}-u_{1,k}^\star\big) \;\ge\; 0,
% \end{aligned}
% \end{equation}
% where $u_{1,k}^\star=-K_k x_k - L_k w_k$ is the optimal control at step $k$.
% \end{theorem}

% \begin{proof}
% We prove optimality via dynamic programming. Define the value function
% \begin{equation}
% \label{eq:value-def}
% V_k(x) \;:=\; \min_{\{u_{1,\ell}\}_{\ell=k}^{N-1}} 
% \left\{ \sum_{\ell=k}^{N-1} \big( x_\ell^\top Q_1 x_\ell + u_{1,\ell}^\top R_1 u_{1,\ell} \big) + x_N^\top Q_{1,N}x_N \right\},
% \end{equation}
% subject to \eqref{eq:general-plant} with the information pattern of Assumption~\ref{ass:measured-wk}. We claim $V_k(x)=x^\top P_k x + c_k$, where $P_k$ is given by \eqref{eq:riccati-Aeff} and $c_k$ is a scalar depending on $\{w_\ell\}_{\ell\ge k}$ but irrelevant for the minimizer.

% \emph{Base case $k=N$.} $V_N(x)=x^\top Q_{1,N}x$, so $P_N=Q_{1,N}$, $c_N=0$.

% \emph{Inductive step.} At time $k$, given $x_k$ and the known $w_k$, the Bellman optimality equation is
% \begin{equation}
% \label{eq:bellman}
% V_k(x_k)=\min_{u_{1,k}}\Big\{ x_k^\top Q_1 x_k + u_{1,k}^\top R_1 u_{1,k} + V_{k+1}(x_{k+1}) \Big\},
% \quad x_{k+1}=A_{\text{eff},k}x_k+B_1u_{1,k}+w_k.
% \end{equation}
% By the induction hypothesis, $V_{k+1}(x)=x^\top P_{k+1}x + c_{k+1}$, so the minimized expression (the ``$Q$-function'') is the strictly convex quadratic
% \begin{equation}
% \label{eq:Qk}
% \mathcal Q_k(x_k,u_{1,k})
% = u_{1,k}^\top H_k u_{1,k}
% + 2\,u_{1,k}^\top B_1^\top P_{k+1}\big(A_{\text{eff},k}x_k+w_k\big)
% + x_k^\top Q_1 x_k 
% + (A_{\text{eff},k}x_k+w_k)^\top P_{k+1}(A_{\text{eff},k}x_k+w_k)
% + c_{k+1},
% \end{equation}
% with $H_k:=R_1+B_1^\top P_{k+1}B_1\succ0$ by $R_1\succ0$.
% Completing the square in $u_{1,k}$ yields the unique minimizer
% \begin{equation}
% \label{eq:u-opt}
% u_{1,k}^\star \;=\; -H_k^{-1}B_1^\top P_{k+1}\big(A_{\text{eff},k}x_k+w_k\big)
% \;=\; -K_k x_k - L_k w_k,
% \end{equation}
% with $K_k,L_k$ in \eqref{eq:KL-opt}. Substituting $u_{1,k}^\star$ into \eqref{eq:bellman} and matching the quadratic term in $x_k$ gives the Riccati update \eqref{eq:riccati-Aeff}. The remaining terms (independent of $u_{1,k}$) define $c_k$, which depends on $w_k$ but does not influence the control law. This proves $V_k(x)=x^\top P_k x + c_k$ and establishes \eqref{eq:KL-opt}.

% \emph{Dominance and gap identity.} For any other choice $\tilde u_{1,k}$,
% \begin{equation}
% \label{eq:one-step-gap}
% \mathcal Q_k(x_k,\tilde u_{1,k}) - \mathcal Q_k(x_k,u_{1,k}^\star)
% = (\tilde u_{1,k}-u_{1,k}^\star)^\top H_k (\tilde u_{1,k}-u_{1,k}^\star) \;\ge\; 0,
% \end{equation}
% by completion of the square and positive definiteness of $H_k$. Summing \eqref{eq:one-step-gap} over $k=0,\dots,N-1$ and using \eqref{eq:bellman} telescopes the value functions, yielding \eqref{eq:gap-sum-of-squares}. In particular, taking $(\tilde K_k,\tilde L_k)=(K^\star_{1,k},0)$ proves \eqref{eq:dominance-stick}.
% \end{proof}


% ++++++++++++++++++++++++++++++++++++++++++++++++++++
% \section{Modified Strategy with Measured Deviations}
% \label{sec:modified-policy}

% We now design Player~1's response when the deviation of Player~2 is \emph{measured at each time step}.
% Let
% \begin{equation}
% w_k := B_2\,\Delta u_{2,k}, \qquad k=0,\dots,N-1,
% \end{equation}
% be the (known) additive input induced by Player~2's deviation.

% \begin{assumption}[Information available to Player~1 at time $k$]
% \label{ass:info-u2}
% At each time $k$, Player~1 knows the current state $x_k$ and the current deviation input $w_k=B_2\Delta u_{2,k}$ (but not future deviations $w_{k+i}$, $i\ge1$). Player~2 continues to apply its FNE feedback $u^\star_{2,k}=-K^\star_{2,k}x_k$.
% \end{assumption}

% Define the ``effective'' open-loop matrix obtained by fixing Player~2 at its FNE law:
% \begin{equation}
% A_{\text{eff},k} := A - B_2 K^\star_{2,k}, \qquad k=0,\dots,N-1.
% \end{equation}
% Then the state evolves as
% \begin{equation}
% x_{k+1} = A_{\text{eff},k} x_k + B_1 u_{1,k} + w_k.
% \end{equation}

% \begin{theorem}[Best response with instantaneous measurement; dominance over stick-with-Nash]
% \label{thm:best-response}
% Under Assumptions~\ref{ass:memoryless-feedback} and \ref{ass:info-u2}, consider the policy class of causal state/input-affine laws
% \begin{equation}
% u_{1,k} = -K_k x_k - L_k w_k, \qquad k=0,\dots,N-1,
% \end{equation}
% where $K_k,L_k$ may be time-varying and depend on model data but not on future $w_{k+i}$.
% Define the Riccati backward recursion
% \begin{equation}
% \label{eq:riccati-eff}
% \begin{aligned}
% &P_N = Q_{1,N}\\
% &\begin{aligned}
%     P_k &= Q_1 + A_{\text{eff},k}^\top P_{k+1} A_{\text{eff},k}
% \\&- A_{\text{eff},k}^\top P_{k+1} B_1 \big(R_1 + B_1^\top P_{k+1} B_1\big)^{-1} B_1^\top P_{k+1} A_{\text{eff},k}.
% \end{aligned}
% \end{aligned}
% \end{equation}
% Then the gains
% \begin{equation}
% \label{eq:KL-best-response}
% \begin{aligned}
% K_k &= \big(R_1 + B_1^\top P_{k+1} B_1\big)^{-1} B_1^\top P_{k+1} A_{\text{eff},k},\\
% L_k &= \big(R_1 + B_1^\top P_{k+1} B_1\big)^{-1} B_1^\top P_{k+1},
% \end{aligned}
% \end{equation}
% minimize Player~1's cost $J_1$ over the stated policy class for any realization of $\{w_k\}$, i.e.,
% \begin{equation}
% \label{eq:dominance-ineq}
% J_1\big(\{K_k,L_k\}\big) \;\le\; J_1\big(\{\tilde K_k,\tilde L_k\}\big)
% \quad \text{for all causal } \{\tilde K_k,\tilde L_k\}.
% \end{equation}
% In particular, the ``stick-with-Nash'' baseline $u_{1,k}=-K^\star_{1,k}x_k$ (which corresponds to $\tilde K_k=K^\star_{1,k},\tilde L_k\equiv 0$) is feasible, hence
% \begin{equation}
% \label{eq:dominance-vs-stick}
% J_1\big(\{K_k,L_k\}\big) \;\le\; J_1\big(\{K^\star_{1,k},0\}\big)
% \quad \text{for every disturbance sequence } \{w_k\}.
% \end{equation}
% \end{theorem}

% \begin{proof}
% Use dynamic programming. At time $k$, given $x_k$ and $w_k$, minimize
% \begin{equation}
% x_k^\top Q_1 x_k + u_{1,k}^\top R_1 u_{1,k} + V_{k+1}(x_{k+1}), \quad x_{k+1}=A_{\text{eff},k}x_k + B_1 u_{1,k} + w_k,
% \end{equation}
% where $V_{k+1}(x)=x^\top P_{k+1}x + c_{k+1}$ with $P_{k+1}$ from \eqref{eq:riccati-eff} (the additive constant $c_{k+1}$ collects known $w$-dependent terms and does not affect the minimizer).
% Expanding and collecting the terms in $u_{1,k}$ gives the strictly convex quadratic
% \begin{equation}
% u_{1,k}^\top \big(R_1 + B_1^\top P_{k+1} B_1\big) u_{1,k}
% + 2\,u_{1,k}^\top B_1^\top P_{k+1} \big(A_{\text{eff},k} x_k + w_k\big) + \text{(terms independent of $u_{1,k}$)}.
% \end{equation}
% Its unique minimizer is
% \begin{equation}
% u_{1,k}^\star \;=\; -\big(R_1 + B_1^\top P_{k+1} B_1\big)^{-1} B_1^\top P_{k+1} \big(A_{\text{eff},k} x_k + w_k\big),
% \end{equation}
% which yields \eqref{eq:KL-best-response}. Substituting $u_{1,k}^\star$ into the Bellman recursion delivers \eqref{eq:riccati-eff}. Since any causal affine law $(\tilde K_k,\tilde L_k)$ is feasible in this DP, optimality implies \eqref{eq:dominance-ineq}, and \eqref{eq:dominance-vs-stick} follows by choosing $(\tilde K_k,\tilde L_k)=(K^\star_{1,k},0)$.
% \end{proof}

% \begin{remark}[Not a new equilibrium]
% The policy in Theorem~\ref{thm:best-response} is Player~1's best response given measured deviations $w_k$ and Player~2 fixed at $K^\star_{2,k}$. The resulting pair is generally \emph{not} a Nash equilibrium.
% \end{remark}

%-----------------------------
% Optional constrained variant
%-----------------------------
% If one wishes to keep Player~1's feedback at the Nash gain and only add feedforward on $w_k$, set
% \begin{equation}
% u_{1,k} = -K^\star_{1,k} x_k - \hat L_k w_k,
% \end{equation}
% and design $\hat L_k$ by running the Riccati recursion on the \emph{homogeneous} closed-loop $A_{\text{cl},k}=A-B_1K^\star_{1,k}-B_2K^\star_{2,k}$:
% \begin{equation}
% \label{eq:riccati-cl}
% \begin{aligned}
% \begin{aligned}
% &\hat P_N = Q_{1,N},\\
% &\begin{aligned}
%     \hat P_k &= S_k + A_{\text{cl},k}^\top \hat P_{k+1} A_{\text{cl},k}\\
%     &- A_{\text{cl},k}^\top \hat P_{k+1} B_1 \big(R_1 + B_1^\top \hat P_{k+1} B_1\big)^{-1} B_1^\top \hat P_{k+1} A_{\text{cl},k},
% \end{aligned}
% \end{aligned}
% \end{aligned}
% \end{equation}
% with $S_k:=Q_1+K_{1,k}^{\star\top}R_1 K^\star_{1,k}$.
% Then
% \begin{equation}
% \label{eq:Lhat}
% \hat L_k = \big(R_1 + B_1^\top \hat P_{k+1} B_1\big)^{-1} B_1^\top \hat P_{k+1}.
% \end{equation}

% \begin{definition}
% \label{def:CF}
% Player~2 applies its FNE law $u_{2,k}=-K^\star_{2,k}x_k$. The compensated controller for Player~1 is
% \[
% u_{1,k}=-K_{1,k}x_k-\hat L_k w_k,
% \]
% where, in general, $K_{1,k}\neq K^\star_{1,k}$ and $\hat L_k$ is the disturbance/preview gain returned by the augmented Riccati recursion.
% \end{definition}

\begin{remark}
\label{rem:CF-homog}
Under the CF policy for Player~1 \eqref{eq:policy-class} and the feedback policy of Player~2 \eqref{eq:feedback nash definition}, the closed-loop state update can be written as
\[
x_{k+1}=\big(A-B_1K_{1,k}-B_2K^\star_{2,k}\big)\,x_k - B_1L_k w_k.
\]
Hence, even for $w_k\equiv 0$, the homogeneous (zero-disturbance) dynamics are
\[
x_{k+1}=\big({A-B_1K_{1,k}-B_2K^\star_{2,k}}\big)\,x_k,
\]
which generally differ from the FNE dynamics
$x_{k+1}=\big(A-B_1K^\star_{1,k}-B_2K^\star_{2,k}\big)\,x_k$.
Stability and performance must therefore be assessed for the Linear Time-Varying (LTV) matrix sequence produced by the augmented Riccati recursion.
\end{remark}


% \begin{proof}
% Direct substitution with $w_k\equiv 0$ gives $x_{k+1}=A_{\text{cl},k}x_k$.
% \end{proof}



% &&&&&&&&&&&&&&&&&&&&&&&&&&&&&&&&&&&&&&&&&&&&&&&&&&&&
% \subsection{Execution-Error Tracking Model and Preview-DBSC Best Response}
% \label{sec:tracking-error-dbsc}

% Even when Player~2 intends to apply its Nash law $u^{\star}_{2,k}=-K^{\star}_{2,k}x_k$, the implemented input $u_{2,k}$ can deviate due to practical effects: actuator/servo dynamics (first-order tracking lag, slew/rate limits, saturations with anti-windup), computation/communication latency and jitter (sample-and-hold delays, packet drops in networked control), inner-loop filtering and smoothing (low-pass prefilters, command shaping), model mismatch and unmodeled cross-couplings, safety/constraint layers that clip or reshape commands (obstacle avoidance, torque/thermal limits), quantization and finite precision, disturbance feedthrough and measurement noise, and setpoint changes that induce transient tracking offsets. These mechanisms produce execution errors $e_k:=u_{2,k}-u^{\star}_{2,k}$ that may be \emph{bias-like} (persistent), \emph{transient/decaying} (lag-induced), or \emph{slowly varying/stochastic} (drift), often correlated across time; the plant experiences $w_k=B_2 e_k$. This motivates simple, identifiable dynamics for $w_k$, such as exponential decay $w_{k+1}=\rho\,w_k$ (inner-loop lag), a random-walk/AR(1) variant $w_{k+1}=\rho\,w_k+v_k$ (drift/noise), or a discrete first-order tracking model for the implemented command that captures both decay and re-injection when $u^{\star}_{2,k}$ changes.

% \paragraph{Implemented-command (first-order lag) dynamics: interpretation and realism.}
% We model Player~2’s actuator/implementation as a discrete-time first-order lag that tracks its desired (Nash) input $u^{\star}_{2,k}=-K^{\star}_{2,k}x_k$. Let $u_{2,k}$ denote the actually implemented command. The lag model
% \begin{equation}\label{eq:lag-time}
% u_{2,k+1}\;=\;\alpha\,u_{2,k}\;+\;(1-\alpha)\,u^{\star}_{2,k},\qquad \alpha\in(0,1),
% \end{equation}
% is the standard sampled counterpart of the continuous-time actuator $\dot u_2(t)=\tfrac{1}{\tau}(u_2^{\star}(t)-u_2(t))$ under zero-order hold with sampling period $\Delta t$, with the identification
% \begin{equation}\label{eq:alpha-tau}
% \alpha \;=\; e^{-\Delta t/\tau}.
% \end{equation}
% This captures two ubiquitous facts: (i) if the target $u^{\star}_{2,k}$ is steady then the tracking error decays geometrically (settling factor $\alpha$), and (ii) changes in $u^{\star}_{2,k}$ inject transient error that then dies out. Such behavior arises from actuator dynamics, command shaping, rate/slew limits, filtering, inner-loop controllers, and latency/smoothing in real systems, and is widely used in control as a realistic low-order approximation of implementation effects.

% % Taking unilateral $z$-transforms of \eqref{eq:lag-time} gives
% % \begin{equation}\label{eq:lag-z}
% % \begin{aligned}
% % z\,U_2(z)\;&=\;\alpha\,U_2(z)\;+\;(1-\alpha)\,U_2^{\star}(z),\\
% % &\Rightarrow
% % \frac{U_2(z)}{U_2^{\star}(z)} \;=\; \frac{1-\alpha}{\,z-\alpha\,}
% % \;=\; \frac{(1-\alpha)\,z^{-1}}{\,1-\alpha z^{-1}\,}.
% % \end{aligned}
% % \end{equation}
% % Thus, the implemented-command discrete transfer function is a stable first-order low-pass with pole at $z=\alpha$ and unit DC gain ($\lim_{z\to 1}U_2/U_2^{\star}=1$). Inverse $z$-transform of \eqref{eq:lag-z} yields the time-domain recursion \eqref{eq:lag-time} directly; equivalently, \eqref{eq:alpha-tau} converts a continuous-time time-constant $\tau$ to the discrete pole $\alpha$.

% % \paragraph{Error/disturbance dynamics (time and $z$ domains).}
% % Define the execution error $e_k:=u_{2,k}-u^{\star}_{2,k}$ and the plant disturbance $w_k:=B_2 e_k$. From \eqref{eq:lag-time},
% % \begin{equation}\label{eq:error-time}
% % e_{k+1}\;=\;\alpha\,e_k\;-\;\Delta u^{\star}_{2,k+1},\qquad
% % \Delta u^{\star}_{2,k+1}:=u^{\star}_{2,k+1}-u^{\star}_{2,k},
% % \end{equation}
% % and, multiplying by $B_2$,
% % \begin{equation}\label{eq:w-time}
% % w_{k+1}\;=\;\alpha\,w_k\;-\;B_2\,\Delta u^{\star}_{2,k+1}
% % \;=\;\alpha\,w_k\;+\;B_2\!\big(K^{\star}_{2,k+1}x_{k+1}-K^{\star}_{2,k}x_k\big).
% % \end{equation}
% % Taking $z$-transforms,
% % \begin{equation}\label{eq:error-z}
% % \begin{aligned}
% % \frac{E(z)}{U_2^{\star}(z)} \;&=\; \frac{U_2(z)-U_2^{\star}(z)}{U_2^{\star}(z)}\\
% % &\;=\;\frac{1-\alpha}{z-\alpha}\;-\;1
% % \;=\;\frac{1-z}{\,z-\alpha\,}\\
% % &\frac{W(z)}{U_2^{\star}(z)} \;=\; B_2\,\frac{1-z}{\,z-\alpha\,}.
% % \end{aligned}
% % \end{equation}
% % Hence, if $u^{\star}_{2,k}$ is slowly varying so that $\Delta u^{\star}_{2,k+1}\!\approx 0$, \eqref{eq:w-time} reduces to the simpler AR(1) model $w_{k+1}\!\approx\!\alpha w_k$; conversely, fast changes in $u^{\star}_{2,k}$ act as a known driving term that re-injects $w_k$.


% We model Player~2's implementation as a first-order lag tracking its Nash command. Let 
% $u^{\star}_{2,k}=-K^{\star}_{2,k}x_k$ denote the desired (Nash) input, $u_{2,k}$ the implemented input, 
% $e_k:=u_{2,k}-u^{\star}_{2,k}$ the execution error, and $w_k:=B_2 e_k$ the disturbance injected at the plant input.

% \begin{assumption}
% \label{ass:lag}
% For some $\alpha\in(0,1)$,
% \begin{equation}\label{eq:u2-lag}
% u_{2,k+1} \;=\; \alpha\,u_{2,k} \;+\; (1-\alpha)\,u^{\star}_{2,k}.
% \end{equation}
% Equivalently,
% \begin{equation}\label{eq:w-rec}
% \begin{aligned}
% w_{k+1} \;&=\; \alpha\,w_k \;-\; B_2\big(u^{\star}_{2,k+1}-u^{\star}_{2,k}\big)\\
% \;&=\; \alpha\,w_k \;+\; B_2\!\big(K^{\star}_{2,k+1}x_{k+1}-K^{\star}_{2,k}x_k\big).
% \end{aligned}
% \end{equation}
% \end{assumption}

% Fix Player~2 at its FNE feedback $u^{\star}_{2,k}=-K^{\star}_{2,k}x_k$ and define the effective matrix
% \begin{equation}\label{eq:Aeff}
% A_{\text{eff},k} \;:=\; A - B_2 K^{\star}_{2,k}.
% \end{equation}
% Player~1 observes the plant
% \begin{equation}\label{eq:plant-lag}
% x_{k+1} \;=\; A_{\text{eff},k}x_k + B_1 u_{1,k} + w_k,
% \end{equation}
% with cost
% \begin{equation}\label{eq:J1-lag}
% \begin{aligned}
% &J_1 \;=\; \sum_{k=0}^{N-1}\Big(x_k^\top Q_1 x_k + u_{1,k}^\top R_1 u_{1,k}\Big) + x_N^\top Q_{1,N}x_N,\\
% &Q_1,Q_{1,N}\succeq 0,\;\; R_1\succ 0.
% \end{aligned}
% \end{equation}

% \begin{lemma}
% \label{lem:aug-lag}
% Let $z_k:=[x_k^\top\; w_k^\top]^\top$. Combining \eqref{eq:w-rec} with \eqref{eq:plant-lag} yields
% \begin{equation}\label{eq:augAB}
% \begin{aligned}
% z_{k+1} &=
% \underbrace{\begin{bmatrix}
% A_{\text{eff},k} & I\\[2pt]
% B_2 K^{\star}_{2,k+1} A_{\text{eff},k} - B_2 K^{\star}_{2,k} & \alpha I + B_2 K^{\star}_{2,k+1}
% \end{bmatrix}}_{\displaystyle \bar A_k}
% z_k\\
% &+
% \underbrace{\begin{bmatrix} B_1\\ B_2 K^{\star}_{2,k+1} B_1\end{bmatrix}}_{\displaystyle \bar B}\,u_{1,k}.
% \end{aligned}
% \end{equation}
% Moreover, \eqref{eq:J1-lag} equals $z_N^\top \bar Q_N z_N + \sum_{k=0}^{N-1}\big(z_k^\top \bar Q z_k + u_{1,k}^\top R_1 u_{1,k}\big)$ with
% \begin{equation}
%     \bar{Q} = \begin{bmatrix} Q_1 & 0 \\ 0 & 0 \end{bmatrix}, \quad \bar{Q}_N = \begin{bmatrix} Q_{1,N} & 0 \\ 0 & 0 \end{bmatrix}.
% \end{equation}
% \end{lemma}

% \begin{theorem}
% \label{thm:lag-optimal}
% Consider causal affine policies for Player~1 of the form
% \begin{equation}\label{eq:policy-lag}
% u_{1,k} \;=\; -K_k x_k - H_k w_k.
% \end{equation}
% Define the finite-horizon Riccati recursion on the augmented system \eqref{eq:augAB}:
% \begin{equation}\label{eq:riccati-aug}
% \begin{aligned}
% &\bar P_N := \bar Q_N,\\
% &\begin{aligned}
% \bar P_k &:= \bar Q + \bar A_k^\top \bar P_{k+1} \bar A_k\\
% &- \bar A_k^\top \bar P_{k+1} \bar B\big(R_1 + \bar B^\top \bar P_{k+1}\bar B\big)^{-1}\bar B^\top \bar P_{k+1} \bar A_k,
% \end{aligned}
% \end{aligned}
% \end{equation}
% with stage curvature $H_k^{\text{(stage)}}:=R_1+\bar B^\top \bar P_{k+1}\bar B \succ 0$. Then the \emph{unique} optimal gains are
% \begin{equation}\label{eq:gains-aug}
% \big[\,K_k\;\;H_k\,\big] \;=\; \big(R_1+\bar B^\top \bar P_{k+1}\bar B\big)^{-1}\,\bar B^\top \bar P_{k+1}\,\bar A_k,
% \end{equation}
% and the corresponding input $u_{1,k}^\star=-K_k x_k - H_k w_k$ minimizes \eqref{eq:J1-lag} over all causal laws of the form \eqref{eq:policy-lag}. Furthermore, for any causal $\{\tilde K_k,\tilde H_k\}$,
% \begin{equation}\label{eq:dom-aug}
% J_1\big(\{\tilde K_k,\tilde H_k\}\big) - J_1\big(\{K_k,H_k\}\big)
% = \sum_{k=0}^{N-1}\big(\tilde u_{1,k}-u_{1,k}^\star\big)^\top \big(R_1+\bar B^\top \bar P_{k+1}\bar B\big)\big(\tilde u_{1,k}-u_{1,k}^\star\big)\;\ge\;0,
% \end{equation}
% where $\tilde u_{1,k}=-\tilde K_k x_k - \tilde H_k w_k$.
% \end{theorem}

% \begin{proof}
% Dynamic programming on $(\bar A_k,\bar B,\bar Q,\bar Q_N,R_1)$ yields a quadratic value function 
% $V_k(z)=z^\top \bar P_k z + \bar c_k$ with $\bar P_k$ given by \eqref{eq:riccati-aug}. 
% At stage $k$, the Bellman minimization is a strictly convex quadratic in $u_{1,k}$ with unique minimizer
% \begin{equation}
% u_{1,k}^\star \;=\; -\big(R_1+\bar B^\top \bar P_{k+1}\bar B\big)^{-1}\bar B^\top \bar P_{k+1}\bar A_k\, z_k,
% \end{equation}
% which, upon partitioning $z_k=[x_k^\top \; w_k^\top]^\top$, gives \eqref{eq:gains-aug}. 
% Completing the square at each stage and summing from $k=0$ to $N-1$ yields the sum-of-squares gap \eqref{eq:dom-aug}.
% \end{proof}





% With the effective plant $x_{k+1}=A_{\mathrm{eff},k}x_k+B_1u_{1,k}+w_k$, the linear recursion \eqref{eq:w-time} yields the augmented system
% \begin{equation}\label{eq:aug-final}
% \begin{bmatrix} x_{k+1}\\ w_{k+1}\end{bmatrix}
% =
% \underbrace{\begin{bmatrix}
% A_{\mathrm{eff},k} & I\\
% B_2 K^{\star}_{2,k+1} A_{\mathrm{eff},k} - B_2 K^{\star}_{2,k} & \alpha I + B_2 K^{\star}_{2,k+1}
% \end{bmatrix}}_{\bar A_k}
% \begin{bmatrix} x_k\\ w_k\end{bmatrix}
% +
% \underbrace{\begin{bmatrix} B_1\\ B_2 K^{\star}_{2,k+1} B_1\end{bmatrix}}_{\bar B} u_{1,k},
% \end{equation}
% on which a standard finite-horizon Riccati recursion produces the preview/DBSC best response
% \begin{equation}\label{eq:dbsc-final}
% u_{1,k}^\star\;=\;-K_k x_k - H_k w_k,\qquad
% \big[\,K_k\;\;H_k\,\big]\;=\;\big(R_1+\bar B^\top \bar P_{k+1}\bar B\big)^{-1}\bar B^\top \bar P_{k+1}\bar A_k,
% \end{equation}
% reducing to the instantaneous-measurement gains when $\alpha\to 0$.


% We refine the execution error model by assuming Player~2 implements its Nash command through a first-order lag. Let
% \begin{equation}
% u^{\star}_{2,k}:=-K^{\star}_{2,k}x_k,\qquad
% e_k:=u_{2,k}-u^{\star}_{2,k},\qquad
% w_k:=B_2 e_k .
% \end{equation}

% \begin{assumption}[Implemented-command lag]
% \label{ass:impl-lag}
% The implemented input follows
% \begin{equation}
% \label{eq:u2-lag}
% u_{2,k+1} \;=\; \alpha\, u_{2,k} \;+\; (1-\alpha)\, u^{\star}_{2,k}, 
% \qquad \alpha\in(0,1),
% \end{equation}
% which implies the error recursion
% \begin{equation}
% \label{eq:e-rec}
% e_{k+1} \;=\; \alpha\, e_k \;-\; \Delta u^{\star}_{2,k+1},\qquad
% \Delta u^{\star}_{2,k+1}:=u^{\star}_{2,k+1}-u^{\star}_{2,k}.
% \end{equation}
% Consequently,
% \begin{equation}
% \label{eq:w-rec}
% w_{k+1} \;=\; \alpha\,w_k \;-\; B_2\,\Delta u^{\star}_{2,k+1}
% \;=\; \alpha\,w_k \;+\; B_2\big(K^{\star}_{2,k+1}x_{k+1}-K^{\star}_{2,k}x_k\big).
% \end{equation}
% \end{assumption}

% Fix Player~2 at its FNE feedback $u^{\star}_{2,k}=-K^{\star}_{2,k}x_k$ and recall
% \begin{equation}
% \label{eq:Aeff-def-again}
% A_{\text{eff},k}:=A - B_2 K^{\star}_{2,k}.
% \end{equation}
% The plant seen by Player~1 is
% \begin{equation}
% \label{eq:plant-seen}
% x_{k+1}=A_{\text{eff},k}x_k + B_1 u_{1,k} + w_k.
% \end{equation}
% Combining \eqref{eq:w-rec} and \eqref{eq:plant-seen} eliminates $x_{k+1}$ and yields the \emph{augmented} linear system
% \begin{equation}
% \label{eq:aug-system}
% \begin{aligned}
%     \begin{bmatrix} x_{k+1}\\ w_{k+1}\end{bmatrix}
% &= 
% \underbrace{\begin{bmatrix}
% A_{\text{eff},k} & I\\[2pt]
% B_2 K^{\star}_{2,k+1} A_{\text{eff},k} - B_2 K^{\star}_{2,k} & \;\alpha I + B_2 K^{\star}_{2,k+1}
% \end{bmatrix}}_{\displaystyle \bar A_k}
% \begin{bmatrix} x_k\\ w_k\end{bmatrix}\\
% & +
% \underbrace{\begin{bmatrix} B_1\\ B_2 K^{\star}_{2,k+1} B_1\end{bmatrix}}_{\displaystyle \bar B}
% \,u_{1,k}.
% \end{aligned}
% \end{equation}
% Player~1's cost is unchanged:
% \begin{equation}
% \label{eq:J1-aug2}
% J_1=\sum_{k=0}^{N-1}\big(x_k^\top Q_1 x_k + u_{1,k}^\top R_1 u_{1,k}\big)+x_N^\top Q_{1,N}x_N,
% \end{equation}
% which we express with block weights $\bar Q:=\mathrm{blkdiag}(Q_1,0)$ and $\bar Q_N:=\mathrm{blkdiag}(Q_{1,N},0)$ on the augmented state $z_k:=[x_k^\top\;w_k^\top]^\top$.

% \begin{theorem}[Best response with implemented-command lag (preview-DBSC)]
% \label{thm:lag-dbsc}
% Under Assumptions~\ref{ass:memoryless-feedback} and \ref{ass:impl-lag}, among causal affine laws
% \begin{equation}
% \label{eq:policy-lag}
% u_{1,k}=-K_k x_k - H_k w_k,
% \end{equation}
% the unique optimal gains are obtained by the Riccati recursion on the augmented system \eqref{eq:aug-system}:
% \begin{equation}
% \label{eq:riccati-aug2}
% \bar P_N:=\bar Q_N,\qquad
% \bar P_k:=\bar Q+\bar A_k^\top \bar P_{k+1}\bar A_k
% -\bar A_k^\top \bar P_{k+1}\bar B\big(R_1+\bar B^\top \bar P_{k+1}\bar B\big)^{-1}\bar B^\top \bar P_{k+1}\bar A_k,
% \end{equation}
% with stage curvature $H_k^{(\mathrm{stage})}:=R_1+\bar B^\top \bar P_{k+1}\bar B\succ0$ and
% \begin{equation}
% \label{eq:KH-aug}
% \big[\,K_k\;\;H_k\,\big]=\big(R_1+\bar B^\top \bar P_{k+1}\bar B\big)^{-1}\bar B^\top \bar P_{k+1}\bar A_k.
% \end{equation}
% Moreover, for any causal $\{\tilde K_k,\tilde H_k\}$,
% \begin{equation}
% \label{eq:dominance-aug}
% J_1\big(\{\tilde K_k,\tilde H_k\}\big)-J_1\big(\{K_k,H_k\}\big)
% =\sum_{k=0}^{N-1}\big(\tilde u_{1,k}-u_{1,k}^\star\big)^\top\! \big(R_1+\bar B^\top \bar P_{k+1}\bar B\big) \big(\tilde u_{1,k}-u_{1,k}^\star\big)\ \ge 0,
% \end{equation}
% where $u_{1,k}^\star=-K_k x_k - H_k w_k$. When $\alpha\to 0$ (very fast implementation), \eqref{eq:aug-system} reduces to the instantaneous-measurement model of Theorem~\ref{thm:best-response}, with $H_k$ collapsing to the instantaneous feedforward gain $L_k$.
% \end{theorem}

% \begin{proof}
% Apply finite-horizon dynamic programming to the augmented linear system \eqref{eq:aug-system} with stage cost $z_k^\top \bar Q z_k + u_{1,k}^\top R_1 u_{1,k}$ and terminal $z_N^\top \bar Q_N z_N$. The value function is quadratic, $V_k(z)=z^\top \bar P_k z + \bar c_k$, with $\bar P_k$ given by \eqref{eq:riccati-aug2}. The Bellman minimizer is
% \begin{equation}
% u_{1,k}^\star=-\big(R_1+\bar B^\top \bar P_{k+1}\bar B\big)^{-1}\bar B^\top \bar P_{k+1}\bar A_k\,z_k,
% \end{equation}
% which yields \eqref{eq:KH-aug} upon partitioning $z_k=[x_k^\top\;w_k^\top]^\top$. The gap identity \eqref{eq:dominance-aug} follows by completion of the square at each stage. The $\alpha\to 0$ reduction is immediate from \eqref{eq:aug-system}.
% \end{proof}

% \begin{remark}[Interpretation]
% Equation \eqref{eq:w-rec} says the error decays as $\alpha^k$ \emph{unless} the desired Nash command changes; changes $\Delta u^{\star}_{2,k+1}$ re-inject error. The augmented dynamics \eqref{eq:aug-system} expose this coupling linearly, allowing a standard preview/DBSC design. All matrices are known given $\{K^{\star}_{2,k}\}$ and the horizon.
% \end{remark}

% ++++++++++++++++++ ORIGINAL +++++++++++++++++++++++

% \begin{theorem}
% \label{prop:OptimalActionP1}
% Consider the game defined in Theorem~\ref{thm:CostAndStateDev} and suppose Player~2 slightly deviates from its feedback Nash strategy. Assume Player~1 can measure or estimate the deviation $\Delta u_{2,k}$ at each time step $k$. Then, Player~1 can optimally compensate for this deviation while maintaining the performance or even improving it by adopting the following modified feedback strategy:
% \begin{equation}
%     \label{eq: optimal strategy for P1 in disturbance situation}
%     u_{1,k} = -K^\star_{1,k}\,x_k - L^\star_k\,B_2\,\Delta u_{2,k},
% \end{equation}
% where the optimal gain matrices \( K^\star_{1,k} \) and \( L^\star_k \) are given by:
% \begin{align}
%     \label{eq: optimal gains}
%     L^\star_k &= \bigl(R_1 + B_1^\top P_{k+1} B_1\bigr)^{-1} B_1^\top P_{k+1}, \\
%     K^\star_{1,k} &= \bigl(R_1 + B_1^\top P_{k+1} B_1\bigr)^{-1} B_1^\top P_{k+1} A_{\mathrm{eff}},
% \end{align}
% with:
% \begin{equation}
%     \label{eq: effective A}
%     A_{\mathrm{eff}} = A - B_2 K^\star_{2,k},
% \end{equation}
% where the matrix \( P_{k+1} \) satisfies the following Difference Riccati Equation (DRE):
% \begin{equation}
%     \label{eq: DRE for P1}
%     \begin{aligned}
%     P_k &= Q_1 + A_{\mathrm{eff}}^\top P_{k+1} A_{\mathrm{eff}} \\
%       & - A_{\mathrm{eff}}^\top P_{k+1} B_1 
%       \bigl(R_1 + B_1^\top P_{k+1} B_1\bigr)^{-1} 
%       B_1^\top P_{k+1} A_{\mathrm{eff}},
% \end{aligned}
% \end{equation}
% with terminal condition $ P_N = Q_{1,N}$.
% \end{theorem}

% \begin{proof}
% Let the effect of Player~2’s deviation be modeled as $w_k := B_2 \Delta u_{2,k}$. Substituting it into the system dynamics (\ref{eq: DT game dynamics}) gives:
% \begin{equation*}
%     x_{k+1} = A_{\mathrm{eff}} x_k + B_1 u_{1,k} + w_k,
% \end{equation*}
% where $A_{\mathrm{eff}}$ is defined in \eqref{eq: effective A}.

% % Adapting the concept introduced in \cite{rugh1996linear}, 
% Inspired by the Disturbance-Based State Compensation (DBSC) approach \cite{rugh1996linear}, which is commonly used in traditional state feedback control design, we assume that the step disturbance $w_k$ is arbitrary but known. We then define the augmented state as
% \begin{equation*}
%     \bar{x}_k = \begin{bmatrix} x_k \\ w_k \end{bmatrix}.
% \end{equation*}

% % Consider $w_k$ is a state of an autonomous dynamical system with unknown initial condition (i.e. $w_{k+1} = ~w_k$, $w_k(0) = w_0$), and 
% % Let $\bar{x}_k$ denote augmented states:

% The state-space representation of the dynamical system can be rewritten in augmented form,
% \begin{equation*}
%     \bar{x}_{k+1} = \bar{A} \bar{x}_k + \bar{B} u_{1,k},
% \end{equation*}
% with
% \begin{equation*}
%     \bar{A} = \begin{bmatrix} A_{\mathrm{eff}} & I \\ 0 & 0 \end{bmatrix}, \quad
%     \bar{B} = \begin{bmatrix} B_1 \\ 0 \end{bmatrix}.
% \end{equation*}

% The objective is to determine optimal control sequence $\{u_{1,k}\}_{k=0}^{N-1}$ that minimizes Player~1's cost function (\ref{eq: LQ cost function}) while mitigating the effect of the disturbance. To this end, we define an adjusted optimization problem for the augmented system as follows:
% \begin{equation}
% \label{eq: augmented LQR}
%     \begin{aligned}
%     &\min_{ \{ u_{1,k} \}_{k=0}^{N-1} } \sum_{k=0}^{N-1} \left( \bar{x}_k^\top \bar{Q} \bar{x}_k + u_{1,k}^\top R_1 u_{1,k} \right) + \bar{x}_N^\top \bar{Q}_N \bar{x}_N, \\
% & \text{subject to:} \quad \bar{x}_{k+1} = \bar{A} \bar{x}_k + \bar{B} u_{1,k}, \quad k =0,\dots, N-1,
% \end{aligned}
% \end{equation}
% where:
% \begin{equation*}
%     \bar{Q} = \begin{bmatrix} Q_1 & 0 \\ 0 & 0 \end{bmatrix}, \quad \bar{Q}_N = \begin{bmatrix} Q_{1,N} & 0 \\ 0 & 0 \end{bmatrix}.
% \end{equation*}

% Equation~(\ref{eq: augmented LQR}) represents a standard finite-horizon discrete-time LQR problem for the augmented system. Its optimal control input is obtained from the augmented state-feedback policy
% \begin{equation*}
%     u_{1,k} = \bar{K}_k \bar{x}_k,
% \end{equation*}
% where the optimal gain $\bar{K}_k$ is computed offline via the recursion \cite{stengel1994optimal},
% \begin{equation}
% \label{eq: optimal augmented gain}
%     \bar{K}_k = (R_1 + \bar{B}^\top \bar{P}_{k+1} \bar{B})^{-1} \bar{B}^\top \bar{P}_{k+1} \bar{A},
% \end{equation}
% with $\bar{P}_{k+1}$ satisfying the following DRE:
% \begin{equation}
%     \label{eq: Riccati for standard LQR}
%     \begin{aligned}
%         \bar{P}_k &= \bar{Q} + \bar{A}^\top \bar{P}_{k+1} \bar{A}  \\&- \bar{A}^\top \bar{P}_{k+1} \bar{B} (R_1 + \bar{B}^\top \bar{P}_{k+1} \bar{B})^{-1} \bar{B}^\top \bar{P}_{k+1} \bar{A}.
%     \end{aligned}
% \end{equation}

% Assume the solution \( \bar{P}_k \) admits the block decomposition:
% \begin{equation*}
%     \bar{P}_k = \begin{bmatrix} P_k^{xx} & P_k^{xw} \\ P_k^{wx} & P_k^{ww} \end{bmatrix}.
% \end{equation*}

% Expanding the optimal gain $\bar{K}_k$ defined in (\ref{eq: optimal augmented gain}) using this block structure for $\bar{P}_{k+1}$, we obtain:
% \begin{align*}
%     \bar{K}_k &= \left(R_1+\begin{bmatrix} B_1 \\ 0 \end{bmatrix}^\top\begin{bmatrix} P_k^{xx} & P_k^{xw} \\ P_k^{wx} & P_k^{ww} \end{bmatrix}\begin{bmatrix} B_1 \\ 0 \end{bmatrix}\right)^{-1}\\
%     & \times \begin{bmatrix} B_1 \\ 0 \end{bmatrix}^\top\begin{bmatrix} P_k^{xx} & P_k^{xw} \\ P_k^{wx} & P_k^{ww} \end{bmatrix} \begin{bmatrix} A_{\mathrm{eff}} & I \\ 0 & 0 \end{bmatrix}.
% \end{align*}

% This expression results in a $2\times 1$ block matrix, which can be partitioned as
% \begin{equation*}
%     \bar{K}_k=\begin{bmatrix} K^\star_{1,k} & L^\star_k \end{bmatrix},
% \end{equation*}
% % \begin{align*}
% %     &\bar{K}_k = \nonumber \\& \begin{bmatrix} (R_1 + B_1^\top P_{k+1}^{xx}B_1)^{-1}B_1^\top P_{k+1}^{xx}A_{\mathrm{eff}} & (R_1 + B_1^\top P_{k+1}^{xx}B_1)^{-1}B_1^\top P_{k+1}^{xx} \end{bmatrix}
% % \end{align*}
% with the optimal gains,
% \begin{align*}
%     K^\star_{1,k} &= (R_1 + B_1^\top P_{k+1}^{xx} B_1)^{-1} B_1^\top P_{k+1}^{xx} A_{\mathrm{eff}}, \\
%     L^\star_k &= (R_1 + B_1^\top P_{k+1}^{xx} B_1)^{-1} B_1^\top P_{k+1}^{xx}.
% \end{align*}

% To compute $P_{k+1}^{xx}$, we substitute the block decomposition $\bar{P}_{k+1}$ into the DRE (\ref{eq: Riccati for standard LQR}). 
% % and expanding it as follows,
% % \begin{align*}
% %     \begin{bmatrix} P_{k}^{xx} & P_{k}^{xw} \\ P_{k}^{wx} & P_{k}^{ww} \end{bmatrix} 
% %     &= \begin{bmatrix} Q_1 & 0 \\ 0 & 0 \end{bmatrix} \notag \\
% %     & + \begin{bmatrix} A_{\mathrm{eff}}^\top & 0 \\ I & 0 \end{bmatrix}
% %     \begin{bmatrix} P_{k+1}^{xx} & P_{k+1}^{xw} \\ P_{k+1}^{wx} & P_{k+1}^{ww} \end{bmatrix}
% %     \begin{bmatrix} A_{\mathrm{eff}} & I \\ 0 & 0 \end{bmatrix} \notag \\
% %     & + \begin{bmatrix} A_{\mathrm{eff}}^\top & 0 \\ I & 0 \end{bmatrix}
% %     \begin{bmatrix} P_{k+1}^{xx} & P_{k+1}^{xw} \\ P_{k+1}^{wx} & P_{k+1}^{ww} \end{bmatrix}
% %     \begin{bmatrix} B_1 \\ 0 \end{bmatrix} \notag \\
% %     & \times \left(R_1 + \begin{bmatrix} B_1^\top & 0 \end{bmatrix}
% %     \begin{bmatrix} P_{k+1}^{xx} & P_{k+1}^{xw} \\ P_{k+1}^{wx} & P_{k+1}^{ww} \end{bmatrix}
% %     \begin{bmatrix} B_1 \\ 0 \end{bmatrix} \right)^{-1} \notag \\
% %     & \times \begin{bmatrix} B_1^\top & 0 \end{bmatrix}
% %     \begin{bmatrix} P_{k+1}^{xx} & P_{k+1}^{xw} \\ P_{k+1}^{wx} & P_{k+1}^{ww} \end{bmatrix} 
% %     \begin{bmatrix} A_{\mathrm{eff}} & I \\ 0 & 0 \end{bmatrix}
% % \end{align*}
% After some algebraic manipulation, the block component $P_{k+1}^{xx}$ satisfies:
% \begin{align*}
%     P_k^{xx} &= Q_1 + A_{\mathrm{eff}}^\top P^{xx}_{k+1} A_{\mathrm{eff}} \nonumber \\&- A_{\mathrm{eff}}^\top P^{xx}_{k+1} B_1 (R_1 + B_1^\top P^{xx}_{k+1} B_1)^{-1} B_1^\top P^{xx}_{k+1} A_{\mathrm{eff}}.
% \end{align*}

% This expression is the standard finite-horizon DRE for the LQR problem applied to the original, non-augmented, disturbance-free system dynamics, with the solution denoted by $P_k$ as in \eqref{eq: DRE for P1}. This concludes the proof.
% \end{proof}



% \begin{proof}
% Defining $A_{\mathrm{eff}}$ as in Equation (\ref{eq: effective A}), considering \(u_{2,k}\) as in Equation \ref{eq: u2 with uncertainty}, and defining $w_k := B_2 \,\Delta u_{2,k}$, we can reduce the problem into a standard LQR problem with step disturbance and rewrite the dynamics defined in Equation~(\ref{eq: DT game dynamics}) as,
%     \begin{equation}
%         \label{eq: modified dynamics}
%         x_{k+1} \;=\; A_{\mathrm{eff}} \, x_k \;+\; B_1\,u_{1,k} \;+\; w_k.
%     \end{equation}

% Adapting the concept introduced in \cite{rugh1996linear}, we assume the step disturbance $w_k$ is arbitrary but perfectly known. Consider $w_k$ is a state of an autonomous dynamical system with unknown initial condition (i.e. $w_{k+1} = ~w_k$, $w_k(0) = w_0$), and let $\bar{x}_k$ denote augmented states:
%     \begin{equation*}
%         \bar{x}_k = \begin{bmatrix} x_k \\ w_k \end{bmatrix},
%     \end{equation*}
    
% Then, the state space representation of the augmented dynamical system can be written as follows,
%     \begin{equation*}
%         \bar{x}_{k+1} = \bar{A} \bar{x}_k + \bar{B} u_{1,k},
%     \end{equation*}
%     where
%     \begin{equation*}
%         \bar{A} = \begin{bmatrix} A_{\mathrm{eff}} & I_{\bar{n}\times \bar{n}} \\ 0 & 0 \end{bmatrix}, \quad
%         \bar{B} = \begin{bmatrix} B_1 \\ 0 \end{bmatrix},  \quad \bar{n} = n_1+n_2
%     \end{equation*}

% Solving the adjusted standard LQR problem with the augmented dynamics as defined as,
%     \begin{equation*}
%         J_1 = \sum_{k=0}^{N-1} \left( \bar{x}_k^\top \bar{Q} \bar{x}_k + u_{1,k}^\top R_1 u_{1,k} \right) + \bar{x}_N^\top \bar{Q}_N \bar{x}_N,
%     \end{equation*}
%     where
%     \begin{equation*}
%         \bar{Q} = \begin{bmatrix} Q_1 & 0 \\ 0 & 0 \end{bmatrix}, \quad \bar{Q}_N = \begin{bmatrix} Q_{1,N} & 0 \\ 0 & 0 \end{bmatrix}, \quad R_1 > 0.
%     \end{equation*}
% returns the optimal policy of,
%     \begin{equation*}
%         u_{1,k} = -\bar{K}_k \bar{x}_k,
%     \end{equation*}
% where
%     \begin{equation}
%         \label{eq: optimal gain in standard LQR}
%         \bar{K}_k = (R_1 + \bar{B}^\top \bar{P}_{k+1} \bar{B})^{-1} \bar{B}^\top \bar{P}_{k+1} \bar{A}.
%     \end{equation}
% and $\bar{P}_{k+1}$ is the solution of the Riccati recursion,
%     \begin{align}
%         \label{eq: Riccati for standard LQR}
%         \bar{P}_k &= \bar{Q} + \bar{A}^\top \bar{P}_{k+1} \bar{A} \nonumber \\&- \bar{A}^\top \bar{P}_{k+1} \bar{B} (R_1 + \bar{B}^\top \bar{P}_{k+1} \bar{B})^{-1} \bar{B}^\top \bar{P}_{k+1} \bar{A},
%     \end{align}
% with the terminal condition $\bar{P}_N = \bar{Q}_N$.

% Expanding the optimal gain obtained in Equation~(\ref{eq: optimal gain in standard LQR}) results in the following expression,
% \begin{align*}
%     \bar{K}_k = \left(R_1+\begin{bmatrix} B_1 \\ 0 \end{bmatrix}^\top\bar{P}_{k+1}\begin{bmatrix} B_1 \\ 0 \end{bmatrix}\right)^{-1}\begin{bmatrix} B_1 \\ 0 \end{bmatrix}^\top\bar{P}_{k+1} \begin{bmatrix} A_{\mathrm{eff}} & I \\ 0 & 0 \end{bmatrix}.
% \end{align*}

% Let assume $P_{k+1}$ is a $2\times 2$ block matrix as follows,
% \[
% \bar{P}_{k+1} = \begin{bmatrix} P_{k+1}^{xx} & P_{k+1}^{xw} \\ P_{k+1}^{wx} & P_{k+1}^{ww} \end{bmatrix}.
% \]

% Substituting the block matrix in the expanded expression results in the following matrix after some manipulations,
% \begin{align*}
%     &\bar{K}_k = \nonumber \\& \begin{bmatrix} (R_1 + B_1^\top P_{k+1}^{xx}B_1)^{-1}B_1^\top P_{k+1}^{xx}A_{\mathrm{eff}} & (R_1 + B_1^\top P_{k+1}^{xx}B_1)^{-1}B_1^\top P_{k+1}^{xx} \end{bmatrix}
% \end{align*}

% Splitting $\bar{K}_k$ into $\bar{K}_k = \begin{bmatrix} K_{1,k}^\star & L_k^\star \end{bmatrix}$ gives the optimal matrices. 

% To obtain the block matrix $P_{k+1}^{xx}$, it is required to expand the Riccati recursion given in Equation~(\ref{eq: Riccati for standard LQR}).
% % \begin{align}
% %     \bar{K}_k &= \left(R_1+\begin{bmatrix} B_1^\top & 0 \end{bmatrix}
% %     \begin{bmatrix} P_{k+1}^{xx} & P_{k+1}^{xw} \\ P_{k+1}^{wx} & P_{k+1}^{ww} \end{bmatrix}
% %     \begin{bmatrix} B_1 \\ 0 \end{bmatrix}\right)^{-1} \notag \\
% %     &\quad \times \begin{bmatrix} B_1^\top & 0 \end{bmatrix}
% %     \begin{bmatrix} P_{k+1}^{xx} & P_{k+1}^{xw} \\ P_{k+1}^{wx} & P_{k+1}^{ww} \end{bmatrix} 
% %     \begin{bmatrix} A_{\mathrm{eff}} & I_{\bar{n}\times \bar{n}} \\ 0 & 0 \end{bmatrix}
% % \end{align}
% Rewriting the Riccati equation,
% \begin{align*}
%         \label{eq: Riccati for standard LQR}
%         \bar{P}_k = \bar{Q} + \bar{A}^\top \bar{P}_{k+1} \bar{A} \nonumber - \bar{A}^\top \bar{P}_{k+1} \bar{B} \bar{K}_k,
%     \end{align*}
% results in,
% \begin{align*}
%     \begin{bmatrix} P_{k}^{xx} & P_{k}^{xw} \\ P_{k}^{wx} & P_{k}^{ww} \end{bmatrix} 
%     &= \begin{bmatrix} Q_1 & 0 \\ 0 & 0 \end{bmatrix} \notag \\
%     & + \begin{bmatrix} A_{\mathrm{eff}}^\top & 0 \\ I & 0 \end{bmatrix}
%     \begin{bmatrix} P_{k+1}^{xx} & P_{k+1}^{xw} \\ P_{k+1}^{wx} & P_{k+1}^{ww} \end{bmatrix}
%     \begin{bmatrix} A_{\mathrm{eff}} & I \\ 0 & 0 \end{bmatrix} \notag \\
%     & + \begin{bmatrix} A_{\mathrm{eff}}^\top & 0 \\ I & 0 \end{bmatrix}
%     \begin{bmatrix} P_{k+1}^{xx} & P_{k+1}^{xw} \\ P_{k+1}^{wx} & P_{k+1}^{ww} \end{bmatrix}
%     \begin{bmatrix} B_1 \\ 0 \end{bmatrix} \notag \\
%     & \times \left(R_1 + \begin{bmatrix} B_1^\top & 0 \end{bmatrix}
%     \begin{bmatrix} P_{k+1}^{xx} & P_{k+1}^{xw} \\ P_{k+1}^{wx} & P_{k+1}^{ww} \end{bmatrix}
%     \begin{bmatrix} B_1 \\ 0 \end{bmatrix} \right)^{-1} \notag \\
%     & \times \begin{bmatrix} B_1^\top & 0 \end{bmatrix}
%     \begin{bmatrix} P_{k+1}^{xx} & P_{k+1}^{xw} \\ P_{k+1}^{wx} & P_{k+1}^{ww} \end{bmatrix} 
%     \begin{bmatrix} A_{\mathrm{eff}} & I \\ 0 & 0 \end{bmatrix}
% \end{align*}

% % \begin{align*}
% %         \begin{bmatrix} P_{k}^{xx} & P_{k}^{xw} \\ P_{k}^{wx} & P_{k}^{ww} \end{bmatrix} & = \begin{bmatrix} Q_1 & 0 \\ 0 & 0 \end{bmatrix} + \begin{bmatrix} A_{\mathrm{eff}}^\top & 0 \\ I & 0 \end{bmatrix}\begin{bmatrix} P_{k+1}^{xx} & P_{k+1}^{xw} \\ P_{k+1}^{wx} & P_{k+1}^{ww} \end{bmatrix}\begin{bmatrix} A_{\mathrm{eff}} & I \\ 0 & 0 \end{bmatrix}\\
% %         & + \begin{bmatrix} A_{\mathrm{eff}}^\top & 0 \\ I & 0 \end{bmatrix}\begin{bmatrix} P_{k+1}^{xx} & P_{k+1}^{xw} \\ P_{k+1}^{wx} & P_{k+1}^{ww} \end{bmatrix}\begin{bmatrix} B_1 \\ 0 \end{bmatrix}\left(R_1+\begin{bmatrix} B_1^\top & 0 \end{bmatrix}\begin{bmatrix} P_{k+1}^{xx} & P_{k+1}^{xw} \\ P_{k+1}^{wx} & P_{k+1}^{ww} \end{bmatrix}\begin{bmatrix} B_1 \\ 0 \end{bmatrix}\right)^{-1}\begin{bmatrix} B_1^\top & 0 \end{bmatrix}\begin{bmatrix} P_{k+1}^{xx} & P_{k+1}^{xw} \\ P_{k+1}^{wx} & P_{k+1}^{ww} \end{bmatrix} \begin{bmatrix} A_{\mathrm{eff}} & I \\ 0 & 0 \end{bmatrix}
% % \end{align*}

% After some manipulation, $P_{k+1}^{xx}$ is obtained from,
% \begin{align*}
%     % P_k^{xx} = Q_1 +  A_{\mathrm{eff}} \\
%     P_k^{xx} &= Q_1 + A_{\mathrm{eff}}^\top P^{xx}_{k+1} A_{\mathrm{eff}} \nonumber \\&- A_{\mathrm{eff}}^\top P^{xx}_{k+1} B_1 (R_1 + B_1^\top P^{xx}_{k+1} B_1)^{-1} B_1^\top P^{xx}_{k+1} A_{\mathrm{eff}},
% \end{align*}
% which is the DRE for non-augmmented non-disturbed system and can be denoted by $P_{k+1}$.

%     % Substituting \( u_k \) into the original system:
%     % \begin{equation}
%     %     x_{k+1} = A x_k + B u_k + w_k.
%     % \end{equation}
%     % \begin{equation}
%     %     x_{k+1} = A x_k + B(-K_k x_k + L_k w_k) + w_k.
%     % \end{equation}
%     % \begin{equation}
%     %     x_{k+1} = (A - B K_k) x_k + (I + B L_k) w_k.
%     % \end{equation}

%     % The disturbance effect is governed by:
%     % \begin{equation}
%     %     \| I + B L_k \|.
%     % \end{equation}

%     % If \( \| I + B L_k \| \ll 1 \), the effect of \( w_k \) is significantly reduced.


    
%     Let the control action of Player~1 be in modified feedback Nash strategy form as defined in Equation (\ref{eq: optimal strategy for P1 in disturbance situation}). This makes an analogy with the standard LQR problem with optimal offline policy addressed in \cite{goel2022power}. The game-theoretic nature of the problem is embedded in $A_{\mathrm{eff}}$ by including Player~2's dynamics.
%     % \hl{However, for the sake of simplicity, the knowledge of future disturbances is not considered in this paper}. 
%     For the linear dynamic system (\ref{eq: modified dynamics}) and the quadratic cost function (\ref{eq: LQ cost function}) for Player~1 ($i=1$), solving the dynamic programming approach \cite{bellman1966dynamic} leads to solving the following \hl{DRE}:
%     \begin{align}
%         P_k &= Q_1 + A_{\mathrm{eff}}^\top P_{k+1} A_{\mathrm{eff}} \nonumber \\
%       & \;-\; A_{\mathrm{eff}}^\top P_{k+1} B_1 
%       \bigl(R_1 + B_1^\top P_{k+1} B_1\bigr)^{-1} 
%       B_1^\top P_{k+1} A_{\mathrm{eff}}, \nonumber \\
%       P_N &= Q_{1,N}.
%     \end{align}
%     The solution of \hl{DRE} gives the optimal control gains (\ref{eq: optimal gains}) for the control policy (\ref{eq: optimal strategy for P1 in disturbance situation}) and ends the proof.
% \end{proof}

% \begin{remark} 
%     The modified feedback policy in Theorem~\ref{prop:OptimalActionP1} introduces an additional feedforward compensation term $i$, allowing Player1 to actively respond to deviations in Player2’s control input. Unlike standard feedback Nash strategies, which react only to the system state, this modification provides direct disturbance rejection, mitigating the adverse effects of Player~2’s deviation before they fully propagate through the system.
% \end{remark}

It is noteworthy that the CF policy is not a Nash equilibrium for the modified interaction, since Player~2 is not best-responding to CF. Joint redesign of both players under implementation dynamics is more involved and is left for future work.

The next section presents a numerical example illustrating the practical utility of Theorems~\ref{thm:lag-optimal}.

% \section{Numerical Illustration}
% \label{sec: numerical example}
% The example problem is adopted from \cite{nortmann2022nash}. We consider a two-player LQ game over a $10$-step horizon. The game dynamics, given by \eqref{eq: DT game dynamics}, is characterized by a discretization step of $\Delta T = 0.4$~s and an initial condition of $x_0 = 1.0$. In the scalar setting, all dynamic matrices are reduced to real-valued scalars. These parameters of scalar dynamics are $A = 1.0$, $B_1 = 1.5$, and $B_2 = -1.0$.

% The weighting scalars in the cost functions given by (\ref{eq: LQ cost function}) are chosen as $Q_1 = Q_{1,N} = 0.1$, $R_1 =0.04$ for Player~1 and $Q_2 = Q_{2,N} =0.2$ for Player~2.

% To illustrate the results of Theorems~\ref{thm:CostAndStateDev} and~\ref{prop:OptimalActionP1}, we numerically analyze and compare three scenarios below.
% \begin{itemize}
%     \item \textbf{Scenario~\uppercase\expandafter{\romannumeral1}}: Both players execute their feedback Nash strategies perfectly.
%     \item \textbf{Scenario~\uppercase\expandafter{\romannumeral2}}: Player~2 slightly deviates from its feedback Nash strategy while Player~1 sticks to its original unadjusted feedback policy. We assumed the control deviation $\Delta u_{2,k}$ at each time step $k$ is a random percentage ranging from $-20\%$ to $+20\%$ of its magnitude $\left|u_{2,k}\right|$. 
%     % The offset is added to avoid the deviation value decay to zero as the magnitude of the control strategy does. 
%     % We have chosen the offset value as $\text{offset}=0.05$.
%     \item \textbf{Scenario~\uppercase\expandafter{\romannumeral3}}: Under the same deviation of Player~2, Player~1 modifies its feedback policy according to (\ref{eq: optimal strategy for P1 in disturbance situation}) from Theorem~\ref{prop:OptimalActionP1} to take advantage of the deviation for a better response.
% \end{itemize}

% The results presented in Table~\ref{tab:costs} clearly illustrate the impact of deviations in Player~2’s feedback Nash strategy on Player~1’s cost. Comparing Scenario~\uppercase\expandafter{\romannumeral1} (feedback Nash) to Scenario~\uppercase\expandafter{\romannumeral2} (feedback), Player~1’s cost nearly doubles, indicating a significant sensitivity to Player~2’s control deviation. However, employing the proposed compensation strategy in Scenario~\uppercase\expandafter{\romannumeral3} (modified feedback) effectively mitigates this adverse effect, reducing Player~1’s cost by approximately $61\%$ compared to the feedback case (Scenario~\uppercase\expandafter{\romannumeral2}). Although Scenario~\uppercase\expandafter{\romannumeral3} still exhibits a $19\%$ increase compared to the feedback Nash, this represents a substantial improvement in robustness relative to the unadjusted feedback case.
% \begin{table}
% \centering
% \caption{Cost Functions Under Different Scenarios}
% \label{tab:costs}
% \renewcommand{\arraystretch}{1.2} % Optional: Adjust row height for better readability
% \definecolor{lightgray}{gray}{0.9} % Adjust as needed
% \begin{tabular}{c >{\columncolor{lightgray}}c c}
% \hline
% \textbf{Scenario} & \textbf{Player~1's Cost} & \textbf{Player~2's Cost} \\ \hline
% (\uppercase\expandafter{\romannumeral1})     & $1.009 \times 10^{-1}$   & $2.06 \times 10^{-1}$\\
% (\uppercase\expandafter{\romannumeral2})     & $3.10 \times 10^{-1}$   & $6.14 \times 10^{-1}$\\
% (\uppercase\expandafter{\romannumeral3})     & $1.207 \times 10^{-1}$   & $2.11 \times 10^{-1}$\\ \hline
% \end{tabular}
% \end{table}

% % \begin{figure}
% %     \centering
% %     \includegraphics[width=1\linewidth]{Figures/u1_plot.png}
% %     \caption{The impact of $\Delta u_{2,k}$ on Player~2's strategy in the three scenarios. \textbf{Top}: Shows the feedback Nash strategy as stated in scenario (\uppercase\expandafter{\romannumeral1}). \textbf{Bottom}: Illustrates the resulting deviation in Player~1's strategy in scenarios (\uppercase\expandafter{\romannumeral2}) and (\uppercase\expandafter{\romannumeral3}).}
% %     \label{fig:u_plot}
% % \end{figure}
% The results in Fig.~\ref{fig:x_plot} illustrate the effectiveness of the modified feedback policy in compensating for the deviation. As shown in the upper plot, the state trajectories under the modified feedback policy (Scenario~\uppercase\expandafter{\romannumeral3}) remain closer to the original feedback Nash trajectory (Scenario~\uppercase\expandafter{\romannumeral1}) compared to the feedback case without compensation (Scenario~\uppercase\expandafter{\romannumeral2}). The lower plot further emphasizes that the state deviations in Scenario~\uppercase\expandafter{\romannumeral3} converge more rapidly to zero, highlighting the effect of the deviation is compensated faster as a result of the feedforward term in the modified feedback policy.
% \begin{figure}
%     \centering
%     \includegraphics[width=1\linewidth]{Figures/xv2.png}
%     \caption{The effect of Player~2's control deviation $\Delta u_{2,k}$ on the game's state for the three scenarios. \textbf{Top}: Compares the joint states in the three scenarios. \textbf{Bottom}: Compares the state deviation of Scenarios \uppercase\expandafter{\romannumeral2} and \uppercase\expandafter{\romannumeral3}.}
%     \label{fig:x_plot}
% \end{figure}

% Fig.~\ref{fig:u_plot} (top) illustrates the effectiveness of the modified feedback policy in compensating for deviations in Player~2's strategy. Player~2's perturbed control input is shown along with Player~1's nominal feedback Nash strategy (Scenario~\uppercase\expandafter{\romannumeral1}) and the modified feedback policy (Scenario~\uppercase\expandafter{\romannumeral3}). The modified feedback control (dash-dot line) closely aligns with the feedback Nash strategy, indicating its capability to mitigate Player~2's deviations effectively.
% The bottom plot explicitly shows Player~1's control deviation across scenarios. The significantly reduced deviation in Scenario\uppercase\expandafter{\romannumeral3}, compared to the unadjusted Scenario~\uppercase\expandafter{\romannumeral2}, further confirms the robustness and improved performance resulting from the proposed policy.
% \begin{figure}
%     \centering
%     \includegraphics[width=1\linewidth]{Figures/allv2.png}
%     \caption{The effect of Player~2's control deviation $\Delta u_{2,k}$ on Player~1's strategy in the three scenarios.}
%     \label{fig:u_plot}
% \end{figure}

% % The simple example above demonstrates the effectiveness of the proposed adjustment to Player~1's strategy when Player~2 deviates from its Nash equilibrium strategy due to execution imperfections.
% It is noteworthy that the proposed policy can effectively leverage Player~2's deviation from Nash equilibrium, resulting in outcomes superior to the original Nash solution. In additional experiments with smaller deviations, the modified policy achieved improvements of up to $0.59\%$ compared to the feedback Nash strategy.

% These findings offer promising opportunities for improving solutions in practical domains such as robotics, traffic management, energy management, etc., where following an exact Nash equilibrium may not always be feasible.

% % % illustrates Player~1’s deviation from its feedback Nash strategy under the gain matrices specified in Theorem~\ref{prop:OptimalActionP1}. The upper plot shows Player~1’s nominal feedback Nash strategy (i.e., in the absence of Player~2’s disturbance), whereas the lower plot compares Player~1’s deviation from feedback Nash strategy with and without policy adjustment. As evidenced by the reduced deviation in the adjusted strategy, revising Player~1’s control policy mitigates the perturbation’s impact. \hl{Unclear. Eventually, is it good to have the strategy very close to the nash or can we achieve a better performance, i.e. lower cost, by employing the adjustment?}

% % Moreover, when Fig.s~\ref{fig:u_plot} and \ref{fig:x_plot} are considered together, it becomes evident that the adjusted policy accelerates state regulation, highlighting both its robustness to disturbance and its efficiency in steering the state to desired values.


% % Notably, the adjusted policy not only preserves, but also improves Player~1's performance beyond the original Nash equilibrium. These results suggest that even under deviation from the Nash equilibrium, Player~1 can exploit the opportunity and a minor adjustment in the feedback strategy restore or even enhance the performance.

% \section{Numerical Example: Two Carts Coupled by a Spring--Damper}
% \label{sec:numerics}

% We consider a two-player, finite-horizon, discrete-time LQ dynamic game with feedback Nash equilibrium (FNE) strategies. Two equal-mass carts are connected by a spring--damper; Player~1 and Player~2 actuate Cart~1 and Cart~2, respectively. Each player regulates its cart to the origin while penalizing control effort. We then impose a realistic execution-lag deviation at Player~2 and show that Player~1 can employ the CF from Theorem~\ref{thm:lag-optimal}, computed via an augmented Riccati recursion, to counteract that structured disturbance while Player~2 remains at FNE.

% Let $x:=\begin{bmatrix}p_1 & v_1 & p_2 & v_2\end{bmatrix}^\top$ in which $p_1$ and $p_2$ are the positions and $v_1$ and $v_2$ are the velocities of Player~1 and~2, respectively. The continuous-time system dynamics are $\dot x = A_c x + B_{1,c} u_1 + B_{2,c} u_2$ with the matrices:
% % \begin{equation}
% % \begin{aligned}
% % \dot p_1 &= v_1,\\
% % m\,\dot v_1 &= -k(p_1-p_2) - c(v_1-v_2) + b_1 u_1,\\
% % \dot p_2 &= v_2,\\
% % m\,\dot v_2 &= \phantom{-}k(p_1-p_2) + c(v_1-v_2) + b_2 u_2,
% % \end{aligned}
% % \end{equation}
% % i.e.,
% \begin{equation*}
% \begin{aligned}
% &A_c = \begin{bmatrix}
% 0&1&0&0\\[2pt]
% -\frac{k}{m}&-\frac{c}{m}&\frac{k}{m}&\frac{c}{m}\\[2pt]
% 0&0&0&1\\[2pt]
% \frac{k}{m}&\frac{c}{m}&-\frac{k}{m}&-\frac{c}{m}
% \end{bmatrix},\\
% &B_{1,c}=\begin{bmatrix}0\\ \frac{1}{m}\\ 0\\ 0\end{bmatrix},\quad
% B_{2,c}=\begin{bmatrix}0\\ 0\\ 0\\ \frac{1}{m}\end{bmatrix}.
% \end{aligned}
% \end{equation*}
% We discretize the dynamics via forward Euler at sampling $\Delta T$. Each player minimizes a quadratic cost \ref{eq:LQ cost function} over a horizon $N$ with the weights $Q_1=\mathrm{diag}(10,\,2,\,0,\,0)$, $ Q_2=\mathrm{diag}(0,\,0,\,10,\,2)$, and $,R_1=R_2=[\,0.05\,]$.
% % \begin{equation}
% % \begin{aligned}
% % x_{k+1} &= A x_k + B_1 u_{1,k} + B_2 u_{2,k},\\
% % A &= I+\Delta T\,A_c,\qquad B_i=\Delta T\,B_{i,c}\ (i=1,2).
% % \end{aligned}
% % \end{equation}

% % \begin{equation}
% % \begin{aligned}
% % &m=0.8,\quad k=0.12,\quad c=30,\quad b_1=b_2=1,\\
% % &\Delta T=0.015~\mathrm{s},\quad N=40,\quad
% % x_0=\begin{bmatrix}0.5&0&-0.5&0\end{bmatrix}^\top.
% % \end{aligned}
% % \end{equation}




% % \begin{equation}
% % \begin{aligned}
% % J_i &= \sum_{k=0}^{N-1}\!\Big(x_k^\top Q_i x_k + u_{i,k}^\top R_i u_{i,k}\Big) + x_N^\top Q_{i,N} x_N,\qquad i\in\{1,2\},\\
% % Q_1&=\mathrm{diag}(10,\,2,\,0,\,0),\quad
% % Q_2=\mathrm{diag}(0,\,0,\,10,\,2),\quad
% % R_1=R_2=[\,0.05\,],\\
% % &Q_{1,N}=Q_{2,N}=Q_1.
% % \end{aligned}
% % \end{equation}
% % The FNE policies $u_{i,k}^\star=-K_{i,k}^\star x_k$ are obtained from the standard coupled backward Riccati recursion for dynamic games, yielding time-varying gains $\{K_{1,k}^\star,K_{2,k}^\star\}_{k=0}^{N-1}$. The induced homogeneous closed loop is
% % \begin{equation}
% % \begin{aligned}
% % A_{\mathrm{cl},k} := A - B_1 K_{1,k}^\star - B_2 K_{2,k}^\star.
% % \end{aligned}
% % \end{equation}

% For error identification purposes, let $u_{2,k}^\star=-K_{2,k}^\star x_k$ denote Player~2’s intended (FNE) command, and let $u_{2,k}$ be the applied input after a first-order discrete lag \eqref{eq:u2-lag} with the $\alpha$ define in Remark~\ref{rem:alpha-tau}. 
% The parameters we considered in this example are given in Table~\ref{tab:params}.

% % \begin{equation}
% % \begin{aligned}
% % y_{k+1} &= \alpha\,y_k + (1-\alpha)\,r_k,\qquad y_0=r_0,\\
% % \alpha &= \exp\!\big(-\Delta T/\tau\big),\qquad \tau=0.8~\mathrm{s}\;\Rightarrow\;\alpha\approx 0.9814.
% % \end{aligned}
% % \end{equation}
% % The plant sees an additive mismatch through the $B_2$-channel,
% % \begin{equation}
% % \begin{aligned}
% % w_k &:= B_2\big(y_k - r_k\big),
% % \end{aligned}
% % \end{equation}
% % so that $u_{2,k}=y_k=r_k+B_2^\dagger w_k$ and the realized state dynamics are
% % \begin{equation}
% % \begin{aligned}
% % x_{k+1} &= \big(A - B_2 K_{2,k}^\star\big) x_k + B_1 u_{1,k} + w_k.
% % \end{aligned}
% % \end{equation}
% % Eliminating $y_k$ and $r_k$ yields the exact augmented dynamics for $z_k:=[x_k^\top\ w_k^\top]^\top$:
% % \begin{equation}
% % \label{eq:exact_augmented_pairs}
% % \begin{aligned}
% % \begin{bmatrix}x_{k+1}\\ w_{k+1}\end{bmatrix}
% % &=
% % \underbrace{\begin{bmatrix}
% % A - B_2K_{2,k}^\star & I\\[4pt]
% % B_2K_{2,k+1}^\star\!\big(A - B_2K_{2,k}^\star\big) - B_2K_{2,k}^\star & \alpha I + B_2K_{2,k+1}^\star
% % \end{bmatrix}}_{\displaystyle \bar A_k}
% % \begin{bmatrix}x_k\\ w_k\end{bmatrix}
% % +
% % \underbrace{\begin{bmatrix}
% % B_1\\[2pt] B_2K_{2,k+1}^\star B_1
% % \end{bmatrix}}_{\displaystyle \bar B_k}u_{1,k}.
% % \end{aligned}
% % \end{equation}
% \begin{table}[h]
% \centering
% \caption{Parameters used in the numerical simulations.}
% \begin{tabular}{lcl}
% \hline
% \multicolumn{3}{c}{\textbf{Model parameters}} \\
% \hline
% $m$       & $1.0$              & Cart mass [kg] \\
% $k$       & $0.12$             & Spring constant [N/m] \\
% $c$       & $30$               & Damping coefficient [N$\cdot$s/m] \\
% $\Delta T$& $0.015$            & Sampling period [s] \\
% $N$       & $40$               & Horizon length [steps] \\
% \hline
% \multicolumn{3}{c}{\textbf{Error/lag parameters}} \\
% \hline
% $\tau$    & $0.8$              & Lag time constant [s] \\
% $\alpha$  & $e^{-\Delta T/\tau}\approx 0.9814$ & Discrete lag factor [–] \\
% \hline
% \multicolumn{3}{c}{\textbf{Initial conditions}} \\
% \hline
% $x_0$     & $\begin{bmatrix}0.5 & 0 & -0.5 & 0\end{bmatrix}^\top$
%           & Initial state [m, m/s, m, m/s] \\
% \hline
% \end{tabular}
% \label{tab:params}
% \end{table}



% \subsection{Test Cases}
% We compare three implementations over the same horizon and initial condition.
% \emph{\textbf{Case~I (FNE)}:} both players apply their finite-horizon feedback Nash laws computed from the coupled Riccati recursion; no actuator lag is present, and Player~2’s intended command is applied directly.
% \emph{\textbf{Case~II (REF)}:} Player~2’s intended command is subjected to a first-order discrete lag with factor $\alpha$, but Player~1 continues to use its nominal feedback policy; this isolates the performance loss due to the execution lag.
% \emph{\textbf{Case~III (CF)}:} the same actuator lag is present on Player~2, while Player~1 designs a compensated feedback by solving a finite-horizon Riccati recursion on the exact augmented realization of the lagged plant, yielding time-varying gains that act on the state and on the lag-induced mismatch signal.


% \begin{figure}
%     \centering
%     \includegraphics[width=1\linewidth]{Figure2/ACC-01 (2).png}
%     \caption{Schematic of the two-cart setup and qualitative final positions under the three cases. 
% Case~I (FNE) shows near-symmetric convergence near the origin. 
% Case~II (REF) illustrates how Player~2’s actuator lag degrades both players’ positions relative to nominal.  
% Case~III (CF) demonstrates how the CF controller enables Player~1 to mitigate the error and somehow recover its nominal outcome, while Player~2 remains misaligned due to its uncompensated lag.}
%     \label{fig:system}
% \end{figure}


% Figure~\ref{fig:state-input comparison} compares the state and input trajectories under the three scenarios. 
% In the nominal FNE (Case~I, solid lines), both players regulate their carts symmetrically to the origin with smooth, coordinated actions. 
% When a first-order execution lag is introduced at Player~2 without compensation (Case~II, dashed lines), the applied input $u_2$ departs significantly from the intended FNE signal, resulting in a sluggish response of Cart~2 ($p_2,v_2$) and, through the spring--damper coupling, a degraded transient for Cart~1 as well. 
% This mismatch explains the pronounced cost increase for both players in Table~\ref{tab:costs}.

% Under the compensated feedback (Case~III, dotted lines), Player~1 augments its state feedback with a disturbance feedforward term tuned to the lag dynamics of Player~2. 
% The effect is visible in $u_1$: the compensated input deviates from the nominal FNE action in order to counteract the actuation shortfall of Player~2. 
% Consequently, the state trajectories $(p_1,v_1,p_2,v_2)$ remain much closer to the nominal FNE benchmark, particularly for Player~1’s variables, confirming that the CF law recovers most of the lag-induced performance loss. 
% Player~2’s cost rises in Case~III, however, since its actual input is still filtered by the lag and it does not adapt its feedback law. 
% This asymmetry reflects the general-sum nature of the game: Player~1 is able to protect its own objective through CF, even if the opponent’s performance worsens.

% Overall, the results illustrate the key insight of Theorem~\ref{thm:lag-optimal}: when a structured deviation such as actuator lag is present at one player, the other can recover its Nash performance by solving an augmented Riccati recursion, without requiring the deviating player to change its policy.


% Numerically (based on Table~\ref{tab:costs}), Player~1’s cost increases by 52\% when the lag is present without compensation (Case~II vs.~Case~I). 
% With CF (Case~III), this penalty is reduced: $J_1$ decreases from $492.40$ to $439.67$, a 10.7\% improvement relative to the uncompensated case. 
% Player~2’s cost, in contrast, rises further under CF since it does not adapt its strategy, underscoring the asymmetric benefit of compensation in a general-sum setting.

% Figure~\ref{fig:system} illustrates the physical setup of the two carts and their final positions under the three cases. 
% In the initial condition, the carts are symmetrically displaced about the origin and connected through the spring--damper. 

% In Case~I (nominal FNE), both players apply their equilibrium feedback laws and drive the carts back toward the origin in a coordinated manner. 
% The final displacements $p_{1,N}^\star$ and $p_{2,N}^\star$ are nearly symmetric and close to zero, reflecting the balanced performance of the Nash solution. 

% In Case~II (lag without compensation), Player~2’s actuator lag prevents it from realizing its intended input. 
% As a result, Cart~2 undershoots relative to its nominal target, and through the coupling the deviation also propagates to Cart~1. 
% The dashed outlines highlight the mismatch from the nominal equilibrium outcome. 

% In Case~III (lag with compensation), Player~1 employs the CF controller. 
% The dotted outlines show that Cart~1’s final displacement is significantly closer to its nominal target than in Case~II, demonstrating recovery of its own performance. 
% Cart~2, however, remains misaligned since it does not alter its strategy and its actuator dynamics still limit the applied force. 

% Overall, the schematic conveys the same conclusion as the numerical trajectories: CF allows the non-deviating player to protect its own Nash performance under structured deviations, even though the deviating player may incur higher cost.



% \begin{figure}
%     \centering
%     \includegraphics[width=1\linewidth]{Figure2/state-input.png}
%     \caption{State trajectories $(p_1,v_1,p_2,v_2)$ and applied inputs under the three cases. The compensated feedback (CF) counters the actuator lag on Player~2 by shaping a mismatch-aware control for Player~1.}
%     \label{fig:state-input comparison}
% \end{figure}


% \begin{table}[h]
% \caption{Finite-horizon costs under the three cases ($\tau=0.8$\,s, $\alpha\approx0.9814$). CF (Case~III) reduces Player~1’s lag penalty relative to Case~II.}
% \centering
% \begin{tabular}{lccc}
% \hline
%  & Case I: FNE & Case II: Lag, no comp. & Case III: CF \\
% \hline
% $J_1$ & 325.31 & 492.40 & \underline{439.67} \\
% $J_2$ & 330.45 & 492.11 & 556.40 \\
% \hline
% \end{tabular}
% \label{tab:costs}
% \end{table}



\section{Numerical Example: Two Carts Coupled by a Spring--Damper}
\label{sec:numerics}

We consider a two-player, finite-horizon, discrete-time LQ dynamic game with feedback Nash equilibrium (FNE) strategies. Two equal-mass carts are connected by a spring--damper; Player~1 and Player~2 actuate Cart~1 and Cart~2, respectively. Each player regulates its cart to the origin while penalizing control effort. We then impose a realistic execution-lag deviation at Player~2 and show that Player~1 can employ the CF from Theorem~\ref{thm:lag-optimal}, computed via an augmented Riccati recursion, to counteract that structured disturbance while Player~2 remains at FNE.
\begin{figure}
    \centering
    \includegraphics[width=1\linewidth]{Figure2/ACC-01.png}
    \caption{Schematic of the two-cart setup and qualitative final positions under the three cases. 
Case~I (FNE) shows near-symmetric convergence near the origin. 
Case~II (REF) illustrates how Player~2’s actuator lag degrades both players’ positions relative to nominal.  
Case~III (CF) shows that the compensated policy enables Player~1 to mitigate the error and approach its nominal outcome, while Player~2 remains misaligned due to its uncompensated lag.}
    \label{fig:system}
\end{figure}
Let $x:=\begin{bmatrix}p_1 & v_1 & p_2 & v_2\end{bmatrix}^\top$, where $p_1$ and $p_2$ are the cart positions, and $v_1$ and $v_2$ are the corresponding  velocities, and $u_1$ and $u_2$ denote the longitudinal actuation forces applied by Player~1 and~2, respectively. 
% \todo[inline]{Reviewer 1, Comment 3}
The continuous-time system dynamics are $\dot x = A_c x + B_{1,c} u_1 + B_{2,c} u_2$ with the matrices:
\begin{equation*}
\begin{aligned}
&A_c = \begin{bmatrix}
0&1&0&0\\[2pt]
-\frac{k}{m}&-\frac{c}{m}&\frac{k}{m}&\frac{c}{m}\\[2pt]
0&0&0&1\\[2pt]
\frac{k}{m}&\frac{c}{m}&-\frac{k}{m}&-\frac{c}{m}
\end{bmatrix},\\
&B_{1,c}=\begin{bmatrix}0\\ \frac{1}{m}\\ 0\\ 0\end{bmatrix},\quad
B_{2,c}=\begin{bmatrix}0\\ 0\\ 0\\ \frac{1}{m}\end{bmatrix}.
\end{aligned}
\end{equation*}
We discretize the dynamics via forward Euler at sampling $\Delta T$. Each player minimizes the quadratic cost in \eqref{eq:LQ cost function} over a horizon $N$ with the weights $Q_1=\mathrm{diag}(10,\,2,\,0,\,0)$, $Q_2=\mathrm{diag}(0,\,0,\,10,\,2)$, and $R_1=R_2=[\,0.05\,]$.

For error identification purposes, let $u_{2,k}^\star=-K_{2,k}^\star x_k$ denote Player~2’s intended (FNE) command, and let $u_{2,k}$ be the applied input after a first-order discrete lag \eqref{eq:u2-lag} with the $\alpha$ defined in Remark~\ref{rem:alpha-tau}. The parameters used are listed in Table~\ref{tab:params}.

\begin{table}[ht]
\centering
\caption{Parameters used in the numerical simulations.}
\begin{tabular}{lcl}
\hline
\multicolumn{3}{c}{\textbf{Model parameters}} \\
\hline
$m$       & $1.0$              & Cart mass [kg] \\
$k$       & $0.12$             & Spring constant [N/m] \\
$c$       & $30$               & Damping coefficient [N$\cdot$s/m] \\
$\Delta T$& $0.015$            & Sampling period [s] \\
$N$       & $40$               & Horizon length [steps] \\
\hline
\multicolumn{3}{c}{\textbf{Error/lag parameters}} \\
\hline
$\tau$    & $0.8$              & Lag time constant [s] \\
$\alpha$  & $e^{-\Delta T/\tau}\approx 0.9814$ & Discrete lag factor [–] \\
\hline
\multicolumn{3}{c}{\textbf{Initial conditions}} \\
\hline
$x_0$     & $\begin{bmatrix}0.5 & 0 & -0.5 & 0\end{bmatrix}^\top$
          & Initial state [m, m/s, m, m/s] \\
\hline
\end{tabular}
\label{tab:params}
\end{table}

\subsection*{Test Cases}
We compare three implementations over the same horizon and initial condition.
\emph{\textbf{Case~I (FNE)}:} both players apply their finite-horizon feedback Nash laws computed from the coupled Riccati recursion; no actuator lag is present, and Player~2’s intended command is applied directly.
\emph{\textbf{Case~II (REF)}:} Player~2’s intended command is subjected to a first-order discrete lag with factor $\alpha$, but Player~1 continues to use its nominal feedback policy; this isolates the performance loss due to the execution lag.
\emph{\textbf{Case~III (CF)}:} the same actuator lag is present on Player~2, while Player~1 designs a compensated feedback by solving a finite-horizon Riccati recursion on the exact augmented realization of the lagged plant, yielding time-varying gains that act on the state and on the lag-induced mismatch signal. Figure~\ref{fig:system} summarizes the two-cart configuration and qualitatively contrasts the resulting final positions under the three cases.

Figure~\ref{fig:state-input comparison} compares the state and input trajectories under the three scenarios. 
In the nominal FNE (Case~I, solid lines), both players regulate their carts symmetrically to the origin with smooth, coordinated actions. 
With a first-order execution lag at Player~2 but no compensation (Case~II, dashed), the applied input $u_2$ deviates from the intended FNE signal (step-like due to the zero-order hold), producing a sluggish response in $(p_2,v_2)$ and, through the spring--damper coupling, a degraded transient in $(p_1,v_1)$. 
This mismatch yields a significant cost increase for both players (Table~\ref{tab:costs}).

\begin{figure}
    \centering
    \includegraphics[width=0.95\linewidth]{Figure2/states_inputs_combo.png}
    \caption{State trajectories $(p_1,v_1,p_2,v_2)$ and applied inputs under the three cases. The compensated feedback (CF) attenuates the effect of Player~2’s actuator lag by shaping a mismatch-aware control for Player~1.}
    \label{fig:state-input comparison}
\end{figure}


Under the compensated feedback (Case~III, dotted), Player~1 modifies its input relative to the nominal FNE action to proactively counteract the actuation shortfall of Player~2. 
The resulting trajectories stay appreciably closer to the FNE benchmark, especially for Player~1, confirming that CF recovers a substantial fraction of the lag-induced performance loss. 
Player~2’s cost, however, rises further because it neither alters its strategy nor removes the actuator lag, highlighting the asymmetric benefit in a general-sum game.


Numerically, Table~\ref{tab:costs} shows that Player~1's cost $J_1$ increases by $51.36\%$ from Case~I to Case~II, when lag is present without compensation. Under CF in Case~III, Player~1's cost decreases to $439.67$, yielding a $10.71\%$ improvement relative to the uncompensated lagged case and recovering about $31.56\%$ of the nominal-to-lag performance loss. At the same time, Player~2's cost $J_2$ increases from $492.11$ in Case~II to $556.40$ in Case~III.

\begin{table}[h]
\caption{Finite-horizon performance under the three cases ($\tau=0.8$\,s, $\alpha\approx0.9814$). CF (Case~III) reduces Player~1’s lag penalty
relative to Case~II.}
\centering
\begin{tabular}{lccc}
\hline
 & Case I: FNE & Case II: Lag, no comp. & Case III: CF \\
\hline
$J_1$        & 325.31   & 492.40   & \textbf{439.67} \\
$J_2$        & 330.45   & 492.11   & 556.40 \\
\hline
$\mathrm{ISE}_1$ & 0.1655   & 0.4566   & {0.2681} \\
$\mathrm{ISE}_2$ & 0.1689   & 0.5076   & {0.3113} \\
\hline
\end{tabular}
\label{tab:costs}
\end{table}

To further assess transient regulation quality, Table~\ref{tab:costs} also reports the finite-horizon integral squared error (ISE) for each player, where $\mathrm{ISE}_1$ corresponds to Player~1 and $\mathrm{ISE}_2$ corresponds to Player~2. These values show the same qualitative improvement for Player~1: relative to Case~II, $\mathrm{ISE}_1$ is reduced by $41.28\%$ under CF. Interestingly, Player~2's $\mathrm{ISE}_2$ also decreases, by $38.67\%$, even though Player~2's cost increases. This indicates that the increase in $J_2$ is not due solely to poorer regulation. Rather, because Player~2's objective penalizes both state deviation and control effort, the modified interaction can reduce regulation error while still increasing the overall quadratic cost.

However, the total cost is increased. This tradeoff is expected in the present formulation. The proposed CF law is obtained by minimizing Player~1's cost over the causal affine class~\eqref{eq:policy-class} while Player~2 remains fixed at its nominal FNE strategy and still suffers from actuator lag. Therefore, the design is intentionally player-specific rather than socially optimal, and a reduction in $J_1$ does not in general imply a reduction in $J_1+J_2$. Mitigating this efficiency tradeoff would require a joint redesign of both players under implementation dynamics, or alternatively a cooperative/team objective, which lies beyond the scope of the present work.


% \todo[inline]{Reviewer 3: Comment 1}


% \begin{table}[h]
% \caption{Finite-horizon costs under the three cases ($\tau=0.8$\,s, $\alpha\approx0.9814$). CF (Case~III) reduces Player~1’s lag penalty relative to Case~II.}
% \centering
% \begin{tabular}{lccc}
% \hline
%  & Case I: FNE & Case II: Lag, no comp. & Case III: CF \\
% \hline
% $J_1$ & 325.31 & 492.40 & \textbf{439.67} \\
% $J_2$ & 330.45 & 492.11 & 556.40 \\
% \hline
% \end{tabular}
% \label{tab:costs}
% \end{table}



\section{Conclusion}
\label{sec: conclusion}
In this paper, we analyzed the impact of control implementation errors in finite-horizon, two-player, nonzero-sum LQ games. 
We quantified how deviations from one player’s nominal
feedback Nash equilibrium input perturb the closed-loop
trajectory and affect the other player’s cost. In particular,
we derived explicit first-order sensitivity expressions that
characterize this dependence in the discrete-time setting.
% \todo[inline]{Before: Our analysis focused on how deviations from the Nash equilibrium influence system stability and performance. We quantified the deviation of the system trajectory from the Nash trajectory and assessed the increase in the cost function when one player fails to follow the exact equilibrium strategy. Our analysis demonstrated that, when a Nash equilibrium exists, the system trajectory remains close to the Nash trajectory for small deviations from the equilibrium strategy.}


Motivated by this sensitivity analysis, we constructed a compensated feedback law for the influenced player by augmenting the nominal dynamics with measurable deviation dynamics. The resulting controller is optimal over a causal affine policy class and includes the uncompensated equilibrium-derived
feedback as a special case.
% \todo[inline]{Before: Additionally, we proposed an adjusted control policy to mitigate the adverse effects of deviations by the other player from the Nash policy. This adjustment, grounded in the principle of dynamic programming, offers superior performance compared to strictly adhering to the pure Nash strategy. By optimally modifying the player’s response, the proposed policy minimizes the negative impact of the other player’s deviations and capitalizes on these deviations to achieve enhanced outcomes.}

A representative numerical example illustrated how the proposed
compensation can recover a substantial portion of the performance
loss induced by actuator lag. These results highlight the value
of exploiting measurable implementation deviations when the
exact execution of nominal Nash policies is not possible.
% \todo[inline]{Before: The effectiveness of the proposed approach was validated through a known example. These findings highlight the potential for improved performance in real-world applications where executing the exact Nash strategy is impractical due to practical challenges such as actuator limitations or external disturbances.}

\subsection*{Future Research Directions}
A promising avenue for future research is the development of online estimation techniques combined with adaptive control laws that can anticipate deviations of the other player from the Nash strategy and dynamically adjust to these errors in real time. Furthermore, examining the effects of more complex deviation patterns, such as time-varying, correlated, or structured errors, could offer deeper insights into the robustness and adaptability of game-theoretic control strategies. 

\section{Acknowledgment}
The authors acknowledge the use of OpenAI’s ChatGPT for structuring content, refining language, and improving clarity. All technical contributions, derivations, and original insights remain the authors' own work.

% \appendix

% \begin{theorem}
%     Optimal Feedforward Gain \( L_k \) for Disturbance Mitigation
%     Consider the discrete-time linear system with a known but random step disturbance $w_k$:
%     \begin{equation}
%         x_{k+1} = A x_k + B u_k + w_k,
%     \end{equation}
%     Define the augmented state:
%     \begin{equation}
%         \bar{x}_k = \begin{bmatrix} x_k \\ w_k \end{bmatrix},
%     \end{equation}
%     leading to the augmented system:
%     \begin{equation}
%         \bar{x}_{k+1} = \bar{A} \bar{x}_k + \bar{B} u_k,
%     \end{equation}
%     where
%     \begin{equation}
%         \bar{A} = \begin{bmatrix} A & I \\ 0 & 0 \end{bmatrix}, \quad
%         \bar{B} = \begin{bmatrix} B \\ 0 \end{bmatrix}.
%     \end{equation}

%     The finite-horizon LQR cost function is given by:
%     \begin{equation}
%         J = \sum_{k=0}^{N-1} \left( \bar{x}_k^\top \bar{Q} \bar{x}_k + u_k^\top R u_k \right) + \bar{x}_N^\top \bar{Q}_N \bar{x}_N,
%     \end{equation}
%     where
%     \begin{equation}
%         \bar{Q} = \begin{bmatrix} Q & 0 \\ 0 & 0 \end{bmatrix}, \quad R > 0.
%     \end{equation}

%     The optimal control policy is:
%     \begin{equation}
%         u_k = -K_k x_k + L_k w_k,
%     \end{equation}
%     where the feedback and feedforward gains are:
%     \begin{align}
%         K_k &= (R + B^\top P_{k+1}^{xx} B)^{-1} B^\top (P_{k+1}^{xx} A + P_{k+1}^{xw}), \\
%         L_k &= (R + B^\top P_{k+1}^{xx} B)^{-1} B^\top (P_{k+1}^{wx} A + P_{k+1}^{ww}).
%     \end{align}

%     The disturbance effect is minimized if:
%     \begin{equation}
%         \| I + B L_k \| \ll 1.
%     \end{equation}

%     The Riccati recursion follows:
%     \begin{equation}
%         P_k = \bar{Q} + \bar{A}^\top P_{k+1} \bar{A} - \bar{A}^\top P_{k+1} \bar{B} (R + \bar{B}^\top P_{k+1} \bar{B})^{-1} \bar{B}^\top P_{k+1} \bar{A},
%     \end{equation}
%     with terminal condition \( P_N = \bar{Q}_N \).

%     The gain \( L_k \) **automatically provides the best possible mitigation** of \( w_k \) for a given horizon length \( N \).
% \end{theorem}

% \begin{proof}
%     We solve the finite-horizon LQR problem for the augmented system. The Riccati recursion satisfies:
%     \begin{equation}
%         P_k = \bar{Q} + \bar{A}^\top P_{k+1} \bar{A} - \bar{A}^\top P_{k+1} \bar{B} (R + \bar{B}^\top P_{k+1} \bar{B})^{-1} \bar{B}^\top P_{k+1} \bar{A}.
%     \end{equation}

%     The optimal control law follows as:
%     \begin{equation}
%         u_k = -\bar{K}_k \bar{x}_k,
%     \end{equation}
%     where
%     \begin{equation}
%         \bar{K}_k = (R + \bar{B}^\top P_{k+1} \bar{B})^{-1} \bar{B}^\top P_{k+1} \bar{A}.
%     \end{equation}

%     Since
%     \begin{equation}
%         \bar{K}_k = \begin{bmatrix} K_k & L_k \end{bmatrix},
%     \end{equation}
%     we extract \( K_k \) and \( L_k \) as:
%     \begin{align}
%         K_k &= (R + B^\top P_{k+1}^{xx} B)^{-1} B^\top (P_{k+1}^{xx} A + P_{k+1}^{xw}), \\
%         L_k &= (R + B^\top P_{k+1}^{xx} B)^{-1} B^\top (P_{k+1}^{wx} A + P_{k+1}^{ww}).
%     \end{align}

%     Substituting \( u_k \) into the original system:
%     \begin{equation}
%         x_{k+1} = A x_k + B u_k + w_k.
%     \end{equation}
%     \begin{equation}
%         x_{k+1} = A x_k + B(-K_k x_k + L_k w_k) + w_k.
%     \end{equation}
%     \begin{equation}
%         x_{k+1} = (A - B K_k) x_k + (I + B L_k) w_k.
%     \end{equation}

%     The disturbance effect is governed by:
%     \begin{equation}
%         \| I + B L_k \|.
%     \end{equation}

%     If \( \| I + B L_k \| \ll 1 \), the effect of \( w_k \) is significantly reduced.

%     Importantly, \( L_k \) is **not arbitrarily chosen** but **derived from the Riccati equation** as part of the optimal solution.

%     \textbf{Conclusion:} The gain \( L_k \) optimally minimizes the impact of \( w_k \), providing a **systematic approach to disturbance rejection within LQR**.
% \end{proof}

% \begin{remark}
%     Unlike manually setting \( L_k = -B^+ \), which ignores the cost function, the proposed method computes \( L_k \) optimally through the Riccati recursion. This ensures that disturbance rejection is balanced with control effort.
% \end{remark}


\bibliographystyle{ieeetr}
\bibliography{references}

\end{document}
